% Template for PLoS
% Version 3.5 March 2018
%
% % % % % % % % % % % % % % % % % % % % % %
%
% -- IMPORTANT NOTE
%
% This template contains comments intended
% to minimize problems and delays during our production
% process. Please follow the template instructions
% whenever possible.
%
% % % % % % % % % % % % % % % % % % % % % % %
%
% Once your paper is accepted for publication,
% PLEASE REMOVE ALL TRACKED CHANGES in this file
% and leave only the final text of your manuscript.
% PLOS recommends the use of latexdiff to track changes during review, as this will help to maintain a clean tex file.
% Visit https://www.ctan.org/pkg/latexdiff?lang=en for info or contact us at latex@plos.org.
%
%
% There are no restrictions on package use within the LaTeX files except that
% no packages listed in the template may be deleted.
%
% Please do not include colors or graphics in the text.
%
% The manuscript LaTeX source should be contained within a single file (do not use \input, \externaldocument, or similar commands).
%
% % % % % % % % % % % % % % % % % % % % % % %
%
% -- FIGURES AND TABLES
%
% Please include tables/figure captions directly after the paragraph where they are first cited in the text.
%
% DO NOT INCLUDE GRAPHICS IN YOUR MANUSCRIPT
% - Figures should be uploaded separately from your manuscript file.
% - Figures generated using LaTeX should be extracted and removed from the PDF before submission.
% - Figures containing multiple panels/subfigures must be combined into one image file before submission.
% For figure citations, please use "Fig" instead of "Figure".
% See http://journals.plos.org/plosone/s/figures for PLOS figure guidelines.
%
% Tables should be cell-based and may not contain:
% - spacing/line breaks within cells to alter layout or alignment
% - do not nest tabular environments (no tabular environments within tabular environments)
% - no graphics or colored text (cell background color/shading OK)
% See http://journals.plos.org/plosone/s/tables for table guidelines.
%
% For tables that exceed the width of the text column, use the adjustwidth environment as illustrated in the example table in text below.
%
% % % % % % % % % % % % % % % % % % % % % % % %
%
% -- EQUATIONS, MATH SYMBOLS, SUBSCRIPTS, AND SUPERSCRIPTS
%
% IMPORTANT
% Below are a few tips to help format your equations and other special characters according to our specifications. For more tips to help reduce the possibility of formatting errors during conversion, please see our LaTeX guidelines at http://journals.plos.org/plosone/s/latex
%
% For inline equations, please be sure to include all portions of an equation in the math environment.
%
% Do not include text that is not math in the math environment.
%
% Please add line breaks to long display equations when possible in order to fit size of the column.
%
% For inline equations, please do not include punctuation (commas, etc) within the math environment unless this is part of the equation.
%
% When adding superscript or subscripts outside of brackets/braces, please group using {}.
%
% Do not use \cal for caligraphic font.  Instead, use \mathcal{}
%
% % % % % % % % % % % % % % % % % % % % % % % %
%
% Please contact latex@plos.org with any questions.
%
% % % % % % % % % % % % % % % % % % % % % % % %

\documentclass[10pt,letterpaper]{article}
\usepackage[top=0.85in,left=2.75in,footskip=0.75in]{geometry}

% amsmath and amssymb packages, useful for mathematical formulas and symbols
\usepackage{amsmath,amssymb}

% Use adjustwidth environment to exceed column width (see example table in text)
\usepackage{changepage}

% Use Unicode characters when possible
\usepackage[utf8x]{inputenc}

% textcomp package and marvosym package for additional characters
\usepackage{textcomp,marvosym}

% cite package, to clean up citations in the main text. Do not remove.
% \usepackage{cite}

% Use nameref to cite supporting information files (see Supporting Information section for more info)
\usepackage{nameref,hyperref}

% line numbers
\usepackage[right]{lineno}

% ligatures disabled
\usepackage{microtype}
\DisableLigatures[f]{encoding = *, family = * }

% color can be used to apply background shading to table cells only
\usepackage[table]{xcolor}

% array package and thick rules for tables
\usepackage{array}

% create "+" rule type for thick vertical lines
\newcolumntype{+}{!{\vrule width 2pt}}

% create \thickcline for thick horizontal lines of variable length
\newlength\savedwidth
\newcommand\thickcline[1]{%
  \noalign{\global\savedwidth\arrayrulewidth\global\arrayrulewidth 2pt}%
  \cline{#1}%
  \noalign{\vskip\arrayrulewidth}%
  \noalign{\global\arrayrulewidth\savedwidth}%
}

% \thickhline command for thick horizontal lines that span the table
\newcommand\thickhline{\noalign{\global\savedwidth\arrayrulewidth\global\arrayrulewidth 2pt}%
\hline
\noalign{\global\arrayrulewidth\savedwidth}}


% Remove comment for double spacing
%\usepackage{setspace}
%\doublespacing

% Text layout
\raggedright
\setlength{\parindent}{0.5cm}
\textwidth 5.25in
\textheight 8.75in

% Bold the 'Figure #' in the caption and separate it from the title/caption with a period
% Captions will be left justified
\usepackage[aboveskip=1pt,labelfont=bf,labelsep=period,justification=raggedright,singlelinecheck=off]{caption}
\renewcommand{\figurename}{Fig}

% Use the PLoS provided BiBTeX style
% \bibliographystyle{plos2015}

% Remove brackets from numbering in List of References
\makeatletter
\renewcommand{\@biblabel}[1]{\quad#1.}
\makeatother



% Header and Footer with logo
\usepackage{lastpage,fancyhdr,graphicx}
\usepackage{epstopdf}
%\pagestyle{myheadings}
\pagestyle{fancy}
\fancyhf{}
%\setlength{\headheight}{27.023pt}
%\lhead{\includegraphics[width=2.0in]{PLOS-submission.eps}}
\rfoot{\thepage/\pageref{LastPage}}
\renewcommand{\headrulewidth}{0pt}
\renewcommand{\footrule}{\hrule height 2pt \vspace{2mm}}
\fancyheadoffset[L]{2.25in}
\fancyfootoffset[L]{2.25in}
\lfoot{\today}

%% Include all macros below

\newcommand{\lorem}{{\bf LOREM}}
\newcommand{\ipsum}{{\bf IPSUM}}


% Pandoc citation processing
\newlength{\csllabelwidth}
\setlength{\csllabelwidth}{3em}
\newlength{\cslhangindent}
\setlength{\cslhangindent}{1.5em}
% for Pandoc 2.8 to 2.10.1
\newenvironment{cslreferences}%
  {}%
  {\par}
% For Pandoc 2.11+
\newenvironment{CSLReferences}[3] % #1 hanging-ident, #2 entry spacing
 {% don't indent paragraphs
  \setlength{\parindent}{0pt}
  % turn on hanging indent if param 1 is 1
  \ifodd #1 \everypar{\setlength{\hangindent}{\cslhangindent}}\ignorespaces\fi
  % set entry spacing
  \ifnum #2 > 0
  \setlength{\parskip}{#2\baselineskip}
  \fi
 }%
 {}
\usepackage{calc} % for calculating minipage widths
\newcommand{\CSLBlock}[1]{#1\hfill\break}
\newcommand{\CSLLeftMargin}[1]{\parbox[t]{\csllabelwidth}{#1}}
\newcommand{\CSLRightInline}[1]{\parbox[t]{\linewidth - \csllabelwidth}{#1}}
\newcommand{\CSLIndent}[1]{\hspace{\cslhangindent}#1}

\usepackage{booktabs}
\usepackage{longtable}
\usepackage{array}
\usepackage{multirow}
\usepackage{wrapfig}
\usepackage{float}
\usepackage{colortbl}
\usepackage{pdflscape}
\usepackage{tabu}
\usepackage{threeparttable}
\usepackage{threeparttablex}
\usepackage[normalem]{ulem}
\usepackage{makecell}
\usepackage{xcolor}
\usepackage{changepage}
\usepackage{booktabs}
\usepackage{longtable}
\usepackage{array}
\usepackage{multirow}
\usepackage{wrapfig}
\usepackage{float}
\usepackage{colortbl}
\usepackage{pdflscape}
\usepackage{tabu}
\usepackage{threeparttable}
\usepackage{threeparttablex}
\usepackage[normalem]{ulem}
\usepackage{makecell}
\usepackage{xcolor}



\usepackage{forarray}
\usepackage{xstring}
\newcommand{\getIndex}[2]{
  \ForEach{,}{\IfEq{#1}{\thislevelitem}{\number\thislevelcount\ExitForEach}{}}{#2}
}

\setcounter{secnumdepth}{0}

\newcommand{\getAff}[1]{
  \getIndex{#1}{Stats,Math,ERHS}
}

\providecommand{\tightlist}{%
  \setlength{\itemsep}{0pt}\setlength{\parskip}{0pt}}

\begin{document}
\vspace*{0.2in}

% Title must be 250 characters or less.
\begin{flushleft}
{\Large
\textbf\newline{Ten simple rules for finding and selecting R
packages} % Please use "sentence case" for title and headings (capitalize only the first word in a title (or heading), the first word in a subtitle (or subheading), and any proper nouns).
}
\newline
% Insert author names, affiliations and corresponding author email (do not include titles, positions, or degrees).
\\
Caroline J. Wendt\textsuperscript{\getAff{Stats}, \getAff{Math}},
G. Brooke Anderson\textsuperscript{\getAff{ERHS}}\textsuperscript{*}\\
\bigskip
\textbf{\getAff{Stats}}Department of Statistics, Colorado State
University, Fort Collins, Colorado, United States of America\\
\textbf{\getAff{Math}}Department of Mathematics, Colorado State
University, Fort Collins, Colorado, United States of America\\
\textbf{\getAff{ERHS}}Department of Environmental \& Radiological Health
Sciences, Colorado State University, Fort Collins, Colorado, United
States of America\\
\bigskip
* Corresponding author: Brooke.Anderson@colostate.edu\\
\end{flushleft}
% Please keep the abstract below 300 words
\section*{Abstract}
R is an increasingly preferred software environment for data analytics
and statistical computing among scientists and practitioners. Packages
markedly extend R's utility and ameliorate inefficient solutions to data
science problems. We outline ten simple rules for finding relevant
packages and determining which package is best for your desired use. We
begin in Rule 1 with tips on how to consider your purpose, which will
guide your search to follow, where, in Rule 2, you'll learn best
practices for finding and collecting options. Rules 3 and 4 will help
you navigate packages' profiles and explore the extent of their online
resources, so that you can be confident in the quality of the package
you choose, and assured that you'll be able to access support. In Rules
5 and 6, you'll become familiar with how the R Community evaluates
packages, and learn how to assess the popularity and utility of packages
for yourself. Rules 7 and 8 will teach you how to investigate and track
package development processes, so you can further evaluate their merit.
We end in Rules 9 and 10 with more hands-on approaches, which involve
digging into package code.

% Please keep the Author Summary between 150 and 200 words
% Use first person. PLOS ONE authors please skip this step.
% Author Summary not valid for PLOS ONE submissions.

\linenumbers

% Use "Eq" instead of "Equation" for equation citations.
\hypertarget{author-summary}{%
\section{Author Summary}\label{author-summary}}

Open source software provides powerful and free tools to facilitate
scientific research, including computational biology. Open source
programming languages like R are created and maintained by international
teams, and a massive community of developers contribute extensions
through open source packages that build on the language's basic
framework. This process has facilitated an expansive set of tools
available to researchers who learn R. However, many researchers new to R
may find it difficult to navigate the available packages, to choose one
that is robust, well-documented, and suited to their task. In this
article, we present ten simple rules to help researchers new to R as
they navigate this vast ecosystem of available R packages.

\hypertarget{disclaimer}{%
\section{Disclaimer}\label{disclaimer}}

G.B. Anderson is a volunteer associate editor at ROpenSci and is an
instructor through the Coursera platform, both of which are mentioned as
potential resources in this article. The views described here reflect
the authors' own views without input from any third party organization.

\hypertarget{funding-acknowledgment}{%
\section{Funding acknowledgment}\label{funding-acknowledgment}}

This research was supported in part by the National Institute of General
Medical Sciences through R25GM132797 (GBA). CJW also received support
from the Honors Undergraduate Program at Colorado State University. The
funders had no role in the manuscript preparation or the decision to
publish.

\hypertarget{introduction}{%
\section{Introduction}\label{introduction}}

As we observe a paradigm shift toward open science {[}1{]}, it's hard
overlook a central axiom surfacing in academic research: computational
reproducibility {[}2,3{]}. The Methods section of scholarly work is
often no longer sufficient to convey data analysis workflows that are
long and complex---data and code must be shared to enable others to
fully understand and replicate many studies. While standards for
disseminating and communicating computational science are evolving with
digitization, responsible researchers are adopting best practices so
others can more easily interpret, validate, and extend their published
knowledge {[}4{]}. These researchers rely on accessible and robust tools
to ensure their computations are reproducible. By using these tools and
implementing other computationally reproducible practices, researchers
themselves directly benefit from improved productivity and
collaboration. In turn, this may reduce some variability in scientific
outcomes that results from differences in workflows; today, the
scientific community deems the absence of automation to be irresponsible
{[}5{]}.

A key tool that upholds the pillars of computational reproducibility is
a scriptable programming language; the R programming language is a
popular choice. This open source language has become a dominant
quantitative programming environment in academic data analysis, enabling
researchers to share workflows and re-execute scripts within and across
subsets of the scientific community {[}5{]}. As the R ecosystem---in
which the life of modern data analysis thrives---rapidly evolves
alongside the burgeoning R community, R is exhibiting sustained growth
when compared to similar languages, especially in academia, healthcare,
and government {[}6{]}. In particular, R is increasingly favored in
computational biology and bioinformatics, two of many disciplines
generating extensive, heterogeneous, and complex data wanting for
heavy-duty data analysis tools that (ideally) support reproducibility
{[}7,8{]}.

R was developed by statisticians, and its base software is
collaboratively maintained by an international group of core
contributors {[}9{]}. The language has an ecosystem of extensions beyond
this base software; unlike several classic proprietary languages such as
MATLAB and SAS, an R user can access, alter, and extend base R software,
then share their work with other users in a formalized way through
extensions called \emph{packages} (as a note, ``package'' is not
perfectly synonymous with ``extension'' because some packages come with
base R and technically aren't extending the base software, but we'll use
``package'' throughout as the commonly recognized colloquial term). An R
package is a fundamental unit of shareable code. There are numerous
analogies in computing between programming and culinary arts: recipe
structures, coding cookbooks, and the like. To conceptualize packages,
imagine you are the chef, base R is the kitchen, and packages are the
special gadgets which allow you to cook and bake new recipes. In
essence, R packages are coding delectables that help you perform
practical tasks and solve problems with interesting techniques.

These user-contributed packages comprise the bulk of enhancements made
to the R environment {[}10{]}. In fact, much of what we know about
statistical methods as well as data analysis and visualization is
wrapped up in R packages---written and documented by R users from around
the world. Are there R packages for wrangling and cleaning data frames,
analyzing gene expression microarray data, or training regression and
classification models? Yes! But how do you \emph{find} an R package to
help you assess the beta diversity of a population, perform
dimensionality reduction, or design interactive applications and
graphics? This answer is not as simple; there are tens of thousands of R
packages. As a natural consequence of the open source nature of R, there
is considerable variation in the quality of R packages and nontrivial
differences among those that provide similar tools.

Packages are essential to venturing beyond base R and, thus, quickly
become an integral aspect of advancing your R capabilities. Those who
have used R packages may know that, although leveraging existing tools
can be advantageous, it can be challenging to find a suitable package
for a given task, which can obstruct potential benefits. An advanced R
user---having developed an intuition for their workflow---may be
relatively confident when searching for and selecting packages. By
contrast, new R users who are unfamiliar with the structure and syntax
of the language may struggle to discover packages because they do not
know where to search, what to look for, or how to sift through options.
When first learning R, two main obstacles are figuring out how to (1)
find a package to accomplish a particular task or solve a problem of
interest and (2) choose the best package to perform that task.

Yet, new users don't need to struggle alone to maneuver these obstacles.
Just as one may go about shopping for shoes, deciding which graduate
program to pursue, or conducting a literature review, there is a science
behind selection. We inform our decisions by assessing, comparing, and
filtering options based on indicators of quality such as utility,
association, and reputation. Likewise, choosing an R package requires
attending to similar details. In this paper, we outline ten simple rules
for finding and selecting R packages, so that you'll spend less time
searching for the right tools and more time coding delightful recipes.

\hypertarget{rule-1-consider-your-purpose}{%
\section{Rule 1: Consider your
purpose}\label{rule-1-consider-your-purpose}}

Usually, there are several ways to accomplish a task while programming,
albeit some more elegant and efficient than others. Before looking for a
package to use for a task, determine whether you need one. Consider your
purpose by first identifying your goal then defining the scope and tasks
to achieve it. For some tasks and workflows, it is practical to code
your own recipe with base R or with packages you're already using, while
other tasks benefit from new tools.

If the scope of your task is simple or reasonable, given your knowledge
and skills, using an R package may not be appropriate. There can be
advantages of coding in base R to complete your task or solve your
problem. First, when you code from scratch, you know precisely what you
are running, and so your script may be easier to decipher and maintain
over time. Conversely, packages require you to rely on shared code with
features or underlying processes of which you may not be aware. Second,
although base R is relatively stable and slow to change, many R packages
evolve rapidly; changes made to packages can cause defects in portions
of your code that can be hard to diagnose.

Packages are favorable in the same sense as kitchen tools: they are
great when a task has a broad or complex scope beyond what you can (or
want to) attempt from scratch. While there are ways to cook up an
algorithm in base R, a relevant package may accomplish your goal with
less code and fewer bugs. Extensive tasks justify sophisticated
frameworks with several functions that form a cohesive package or even a
suite of related packages. Data manipulation is one such common task
that has been streamlined by packages such as \texttt{dplyr} and
\texttt{tidyr} {[}11,12{]}, both part of the \texttt{tidyverse} suite
(Table \ref{table1}). Nevertheless, don't be discouraged if you have a
task that seems too unusual for a package. There are indeed packages for
seemingly singular tasks, which you may favor over coding from scratch.
You can use an R package to access the Twitter API, send emails from R,
and there is even a package for converting English letters to numbers as
on a telephone keypad {[}13--15{]}.

\begin{table}[!t]
\begin{adjustwidth}{-2.25in}{0in} % Comment out/remove adjustwidth environment if table fits in text column.
\caption{Common R packages for general purposes}
\centering
\resizebox{\linewidth}{!}{
\begin{tabular}[t]{>{\raggedright\arraybackslash}p{10em}>{\raggedright\arraybackslash}p{20em}lll}
\toprule
\textbf{Package} & \textbf{Description} & \textbf{Year} & \textbf{Author} & \textbf{Documentation}\\
\midrule
\addlinespace[0.3em]
\multicolumn{5}{l}{\textbf{Data Manipulation}}\\
\hspace{1em}readr & Read rectangular data (e.g., csv, tsv, and fwf) & 2015 & Wickham et al. \cite{@readr} & https://readr.tidyverse.org/\\
\hspace{1em}dplyr & Grammar of tidy data manipulation & 2014 & Wickham et al. \cite{@dplyr} & https://dplyr.tidyverse.org/\\
\hspace{1em}tidyr & Tidy messy data & 2014 & Wickham \& Henry \cite{@tidyr} & https://tidyr.tidyverse.org/\\
\addlinespace[0.3em]
\multicolumn{5}{l}{\textbf{Statistical Modeling}}\\
\hspace{1em}broom\textsuperscript{*} & Tidy model output; convert statistical objects into tidy tibbles & 2014 & Robinson \& Hayes \cite{@broom} & https://broom.tidymodels.org/\\
\hspace{1em}purrr & Functional programming tools for functions and vectors & 2015 & Henry \& Wickham \cite{@purrr} & https://purrr.tidyverse.org/\\
\hspace{1em}caret & Framework for predictive modeling & 2007 & Kuhn \cite{@caert} & https://topepo.github.io/caret/index.html\\
\addlinespace[0.3em]
\multicolumn{5}{l}{\textbf{Data Visualization}}\\
\hspace{1em}ggplot2 & System for data visualization & 2007 & Wickham et al. \cite{@ggplot2} & https://ggplot2.tidyverse.org/\\
\hspace{1em}kableExtra & Build complex tables and manipulate styles & 2017 & Zhu \cite{@kableExtra} & https://haozhu233.github.io/kableExtra/\\
\hspace{1em}rmarkdown & Authoring framework for data science and reproducible research & 2014 & Allaire et al. \cite{@rmarkdown} & https://rmarkdown.rstudio.com/lesson-1.html\\
\bottomrule
\multicolumn{5}{l}{\textsuperscript{*} See the biobroom \cite{@biobroom} analog in Bioconductor}\\
\end{tabular}}
\label{table1}
\end{adjustwidth}
\end{table}

If you decide you need a package, next ask yourself: Which
functionalities in base R are restrictive in context of your task? Which
new functionalities would expand what you can do? If you identify
limitations of your current toolbox before searching for new tools,
you'll be primed to recognize what you do and don't need. List
domain-specific keywords that describe what you are trying to do and how
you could do it, so you narrow your search. Identify the type of inputs
you have and envision working with them; contemplate the desired outputs
and corresponding format. Suppose you are using Bioconductor packages in
analyses and have outputs you want to visualize; you must consider that
the inputs may be of a certain class---a type of S4 object, for example
{[}16{]}. When you look for a package to visualize the data, you'll want
to choose one that harmoniously handles this type of object as an input.
If you don't quite know (or have a difficult time articulating) what you
need, don't hesitate to search the internet for examples that capture
your purpose. It can be helpful to search with Google (and Google Images
for plots and figures) to identify the terminology that describes what
you are trying to do, how people approach similar tasks, and potential
challenges that may arise.

\hypertarget{rule-2-find-and-collect-options}{%
\section{Rule 2: Find and collect
options}\label{rule-2-find-and-collect-options}}

You'll encounter tips and leads for finding R packages in various places
online, in print, and elsewhere. If you've searched for packages a lot,
you likely have a shortlist of tried and true starting points---a perk
of being a long-time R user. On the other hand, new R users haven't yet
developed that internal compass, and they often don't know where to
start looking for a package to suit their task. Here are a few
directions to head in: You can discover new packages any time you learn
R-related topics, browse the internet, or connect with the international
community of R users.

When learning how to program in R, you are typically introduced to some
of the most common packages, which tend to have more general purposes.
Table \ref{table1} gives some examples of popular packages you may have
learned about in an introductory R course or book, particularly if it
focused on the \texttt{tidyverse} approach to coding in R. In addition,
reputable online tutorials, courses, and books are helpful resources for
acquiring knowledge about packages that are versatile and
reliable---they tend to be concise, accessible, and either affordable or
offered at \href{https://committedtotape.shinyapps.io/freeR/}{no cost}
to the learner. For example,
\href{https://www.codecademy.com/learn/learn-r}{Codeacademy},
\href{https://www.coursera.org/learn/r-programming}{Coursera}, while
other Massive Open Online Courses provide platforms for interactive
learning, and numerous R programming books are available freely
online---many of which are highlighted on the
\href{https://bookdown.org/}{bookdown.org site} or obtainable through
\href{https://rstudio.com/resources/books/}{RStudio's books}---or in
print series like the Springer UseR! Series and the Chapman \& Hall/CRC
R Series.

You can also browse online to look for R packages. Search engines such
as Google return ample pages related to anything
\texttt{"\ldots{}\ in\ R"}. This can be a good place to start,
especially if you use the list of terms you developed for Rule 1. For a
more directed approach, you can search the package lists on repositories
including the \href{https://cran.r-project.org/}{Comprehensive R Archive
Network (CRAN)}, \href{https://www.bioconductor.org/}{Bioconductor}, and
\href{https://ropensci.org/}{rOpenSci} or browse packages available on
code-sharing websites like \href{https://github.com/}{GitHub} and
\href{https://about.gitlab.com/}{GitLab} (see Rule 3), with language
filters for finding repositories that heavily rely on R code.

You can also search with a more targeted approach by starting from a
curated list of R packages on a particular topic.
\href{https://cran.r-project.org/web/views/}{CRAN Task Views}
concentrate on topics from certain disciplines or methodological
approaches (e.g., Econometrics, Genetics, Optimization, Spatial). In the
HTML version, you can browse alphabetized subcategories of packages
within each Task View and read concise descriptions to find tools with
specific functionality (Fig \ref{fig1}). Alternatively, you can access
Task Views directly from the R console with \texttt{ctv::CRAN.views()}
{[}17{]} (note that here and throughout, we use this compact
\texttt{package::function()} notation to indicate a function from a
specific package). There are currently 41 Task Views that collectively
contain thousands of packages, all curated and regularly tested.
Moreover, CRAN Task Views provide tools that enable you to automatically
install all packages within a targeted area of interest. Ultimately, by
providing task-based organization, easy simultaneous installation of
related packages, meta-information, ensured maintenance, and quality
control, CRAN Task Views address several major user-end issues that have
arisen due to the sheer quantity of available packages {[}17{]}.

\begin{figure}[!h]
\caption{{\bf A snapshot of CRAN Task Views.}
CRAN Task Views offer listings of available R packages within specific
topic areas, like time series, optimization, medical imaging, genetics, 
and high-performance computing.}
\label{fig1}
\end{figure}

The computational biology and bioinformatics analogs of CRAN Task Views,
\href{https://www.bioconductor.org/packages/release/BiocViews.html\#___Software}{biocViews},
are lists of packages available through the Bioconductor repository that
provide tools for analyzing biological data (Fig \ref{fig2}). Your
keywords from Rule 1 will come in handy here, too; for instance, you can
search ``RNA-seq'' to find all packages related to high-throughput
sequencing methods for analyzing RNA populations. You can also browse
the biocViews tree menu to narrow your search based on which general and
specific categories describe your task. The list includes each package's
name linked to information and documentation, its maintainer, a short
description, and its rank relative to other Bioconductor packages.

\begin{figure}[!h]
\caption{{\bf A snapshot of biocViews.}
Bioconductor---which serves as a repository for packages focused on
high-throughput genomic data---offers similar curated lists of packages through
biocViews, including listings by assay domain, by biological question, and by
statistical method.}
\label{fig2}
\end{figure}

Other curated collections and topically-linked lists of packages are
also available.
\href{http://dirk.eddelbuettel.com/cranberries/index.html}{CRANberries},
for example, is a hub of information about new, updated, and removed
packages from CRAN. Another place to find well-maintained tools is
\href{https://ropensci.org/packages/}{rOpenSci packages}. These
filterable and searchable peer-reviewed (see Rule 6) R packages are
organized by name, maintainer, description, and status markers based on
activity, association, and review (Fig \ref{fig3}).

\begin{figure}[!h]
\caption{{\bf An example of the rOpenSci package search page, filtered to image processing packages.}
rOpenSci is a non-profit initiative committed to the advancement of open science practices through the development, maintenance, and promotion of peer-reviewed R packages, which are easily searchable on their webpage.}
\label{fig3}
\end{figure}

You'll also find tips on R packages by joining in R's vibrant
international community. An inclusive and collaborative community is an
overlooked, yet integral, aspect of a software's success {[}18{]}. The R
community has a widespread internet presence across various platforms;
however, members are markedly active on Twitter, a place where R users
seek help, share ideas, and stay informed on \texttt{\#rstats}
happenings, including releases of new packages {[}19{]}. R-related blogs
serve as another informal, up-to-date, and more detailed avenue for
communicating and promoting R-related information. For example, Joseph
Rickert, Ambassador at Large for RStudio, writes monthly posts on the
\href{https://rviews.rstudio.com/}{R Views} blog highlighting
interesting new R packages. Rickert also features special articles about
recently released packages and lists of top packages within certain
categories, including: Computational Methods, Data, Machine Learning,
Medicine, Science, Statistics, Time Series, Utilities, and
Visualization.

Members of the R community gather regularly at conferences to share
their latest work. You can attend---or access content from---these
conferences to learn about improved, recently released, or up-and-coming
R package developments and applications. Two large annual R conferences
with broad scopes are
\href{https://rstudio.com/conference/}{rstudio::conf} and
\href{https://www.r-consortium.org/}{useR!}. There are also a number of
smaller, more discipline-tailored R conferences, including:
\href{https://rinpharma.com/}{R/Pharma} (pharmaceutical development),
\href{https://www.rinfinance.com/}{R/Finance} (applied finance), and
\href{https://www.bioconductor.org/help/events/}{BioC} (open source
software for bioinformatics). Conferences in your field may foster
connections with fellow scientists who use R for similar tasks and help
you collect information about packages related to your expertise. Talks
and presentations at conferences are often recorded and made available
online for playback if you can't attend in person or wish to revisit
content from past conferences. For instance, R Studio has a
\href{https://rstudio.com/resources/rstudioconf-2020/}{webpage} devoted
to videos of past conferences by year and the talks from useR! are
posted on the R Consortium
\href{https://www.youtube.com/channel/UC_R5smHVXRYGhZYDJsnXTwg/featured}{YouTube
channel}. The R community and rOpenSci have also popularized
nontraditional Unconferences and R Collaboratives at which people
interact and share ideas in a collegial and unstructured atmosphere; as
you may imagine, R packages are a hot topic of conversation at these
events.

Don't forget, a developer of an R package may intend for it to be
private (exclusively for personal or professional use) or public (at no
cost and available for use by anyone) {[}20{]}. If your task is specific
to a line of research, consult colleagues to see if they have a relevant
(private) package they'd be willing to share.

\hypertarget{rule-3-check-how-its-shared}{%
\section{Rule 3: Check how it's
shared}\label{rule-3-check-how-its-shared}}

R developers share their packages through several platforms. The primary
type of platform is a repository, where developers share packages for
public use. An R package repository is essentially a warehouse for tools
before they are shipped to your kitchen; in computing terms, a
repository provides a central online location in which packages are
stored and managed. Some repositories provide vetting mechanisms that
tame unwieldy aspects of the R ecosystem by regularly checking
underlying code and managing corresponding webs of dependencies. Two
traditional repositories for R packages are CRAN and Bioconductor;
however, there are lesser-known repositories to explore. Developers can
also share their R packages through large code-sharing sites, including
public version control platforms like GitHub and GitLab; these have
fewer restrictions on the format or content of shared code compared to
repositories. For more private sharing, rather than making packages
accessible to everyone via internet repositories, some developers share
their code in zipped files directly with collaborators (see Rule 2).

The CRAN repository is the most established R package repository, and it
serves as the primary source from which you'll likely install R
packages. The Institute for Statistics and Mathematics of
Wirtschaftsuniversität Wien hosts the CRAN server and Contributors from
the R Core Team (established in 1997) maintain its monumental collection
of R packages (over 16,000 as of August 2020) on almost any conceivable
topic, from random forests to multidimensional maps {[}9,21,22{]}. Due
to its longevity and historical role, Rickert asserts that ``CRAN is the
greatest open source repository for statistical computing knowledge in
the world'' {[}20{]}. A key advantage of using a package from CRAN is
that the repository integrates easily with your base R installation
{[}23{]}. You can simply use the \texttt{install.packages()} function
with default values to install a package from CRAN, along with all of
the packages it depends on; source and/or binary code are automatically
saved to your computer in a designated package library {[}9,24{]}. When
you want to use a particular package, you load it to your R session via
the \texttt{library()} function from base R {[}9{]}. These features
ensure that integration of a new package with your version of the base R
software is typically silent and seamless.

CRAN imposes some regulations on the submission and maintenance of its
hosted packages. For example, a package must pass a series of initial
stability tests in accordance with the CRAN Repository Policy before
obtaining publication privileges {[}25{]}. CRAN does not, however, have
a peer-review process (see Rule 6), and so these checks do not extend to
the substance of the package, in terms of ensuring it works as
advertised in the documentation. At worst, then, a CRAN package might
have a solid framework with help files for each outward-facing function,
but its underlying code might not necessarily do what it claims. CRAN
maintainers actively monitor contributed packages to ensure
compatibility with the latest version of R and dependencies; they update
metadata accordingly and ask for modifications or remove packages that
don't uphold these publication quality standards. For instance, to see
the metadata page for the \texttt{ggplot2} package, you can go to
\url{https://CRAN.R-project.org/package=ggplot2}; replace the package
name at the end of the web address with the name of any other CRAN
package to access analogous package-specific metadata (Fig \ref{fig4})
{[}26{]}.

\begin{figure}[!h]
\caption{{\bf An example of the CRAN page for the \texttt{ggplot2} package.}
Each package on CRAN has its own webpage, which can be used to follow several of the Rules highlighted in this article. This snapshot of the page for the 
\texttt{ggplot2} package indicates where information related to numerous Rules can be found on the package's CRAN page.}
\label{fig4}
\end{figure}

If you're looking for tools specifically designed for high-throughput
genomic data, the Bioconductor repository is a more specialized
repository than CRAN that's worth investigating {[}27{]}. The
Bioconductor project was motivated by a need for transparent,
reproducible, and efficient software in computational biology and
bioinformatics; it supports the integration of computational rigor and
reproducibility in research on biological processes {[}7{]}.
Bioconductor packages facilitate the analysis and comprehension of
biological data and help users solve problems that arise when working
with high-throughput genomic data such as those related to microarrays,
sequencing, flow cytometry, mass spectrometry, and image analysis
{[}27{]}. Bioconductor boasts a modularized design, wherein data
structures, functions, and the packages that contain them have distinct
roles that are accompanied by comprehensive documentation. For example,
the page for the package \texttt{DESeq2} can be found at
\url{https://bioconductor.org/packages/release/bioc/html/DESeq2.html};
you can view pages corresponding to other Bioconductor packages by
altering the web address as previously described (Fig \ref{fig5})
{[}28{]}. Accessibility is a pillar of the Bioconductor project; all
forms of documentation, including courses, vignettes, and interactive
documents, are curated for individuals with expertise in adjacent
disciplines and minimal experience with R (see Rule 4) {[}7{]}.
Bioconductor packages also integrate well with base R by different means
than CRAN packages; you should use \texttt{BiocManager::install()} to
install Bioconductor packages with their corresponding dependencies
{[}29{]}.

\begin{figure}[!h]
\caption{{\bf An example of the Bioconductor page for the \texttt{DESeq2} package}
Each package on Bioconductor has its own webpage, which can be used to follow many Rules discussed in this article. This snapshot of the page for the \texttt{DESeq2} package indicates where information related to several of the Rules can be found on the package's Bioconductor page.}
\label{fig5}
\end{figure}

Bioconductor has a structured release schedule and thus a unique package
management system that automatically loads your Bioconductor packages in
versions that are all from the same release cycle and compatible with
your current version of R {[}27{]}. This helps, among other things, keep
metadata properly aligned with sample data. Bioconductor has strict
criteria for package submissions: a package must be relevant to
high-throughput genomic analysis, interoperable with other Bioconductor
packages, well-documented, supported in the long-term, and comply with
additional package guidelines {[}30{]}. To ensure Bioconductor software
maintains a centralized content focus, their packages are coordinated
and peer-reviewed. The peer-review processes includes a
technically-focused review by a primary reviewer and may also include
secondary reviews by content-matter experts {[}27{]}.

In addition to CRAN and Bioconductor, there are other smaller, more
specialized repositories that share R packages. The non-profit
organization, rOpenSci, for example, runs a repository as part of their
commitment to promote open science practices through technical and
social infrastructure for the R community {[}31{]}. The repository only
includes packages that have either been developed by staff to fill a gap
in infrastructure or have passed their open review process. These
packages are within the scope and aims defined by rOpenSci, including
those that focus on data extraction, workflow automation, and field or
laboratory reproducibility tools {[}32{]}. At a smaller scale, R package
maintainers can create and host their own personalized repository using
the \texttt{drat} (Drat R Archive Template) package, which enables
developers to design individual repositories and suites of coordinated
repositories for packages that are stored in and/or distributed through
platforms like GitHub {[}33{]}. In some cases, this might be done to
host a package that doesn't meet certain restrictions from other
repositories (e.g., to share data through a package that's larger than
CRAN's size limit for packages {[}34{]}). There are other platforms
strictly for the development, rather than distribution, of R packages
such as R-Forge and Omegahat, which are beyond the scope of this paper
{[}35,36{]}.

Packages can also be hosted on online code sharing sites like GitHub or
GitLab, both of which are online hosting platforms for repositories
tracked via git version control. While there are evident benefits of
exclusive R package repositories, there are also considerable advantages
to hosting a package on these online platforms {[}20,37{]}. First,
package developers can post the directory with all source code for their
package as they develop it, from the earliest stages, rather than
waiting until the package is publication-ready. What's more, these
platforms allow for continuous integration and continuous deployment for
testing data and code to ensure a stable end-product.

One particularly popular and effective platform for sharing R packages
is GitHub (Fig \ref{fig6}). An increasing number of packages are hosted
on GitHub during the development stages---it will often be the first
place you can find up-and-coming packages (see Rule 2). If stable
release versions of these packages are published on CRAN or Bioconductor
{[}38{]}, GitHub is still a good place to explore planned changes for
the next version that the developer will soon release. GitHub provides R
users with open access to package code, a timeline of help resources
(see Rule 4), and a direct line of communication to developers. Both R
users and package developers benefit from interactive feedback channels
through GitHub Issues and the Star rating system.

\begin{figure}[!h]
\caption{{\bf An example of the GitHub page for the \texttt{bookdown} package.}
Developers utilize the GitHub platform to share their packages in real-time for public use and feedback. This snapshot of the page for the \texttt{bookdown} package labels where information related to several of the Rules can be found on the package's GitHub page.}
\label{fig6}
\end{figure}

It's straightforward to install the latest version of a package you want
to use from GitHub via
\texttt{devtools::install\_github("developer/package")}. However, the
decentralized nature of GitHub isn't conducive to a tool that
automatically locates and installs corresponding dependencies that are
not on CRAN {[}39{]}. In fact, the rapid uptick in package development
and ensuing inter-repository dependencies has sparked an ongoing debate
on whether regulated repositories such as CRAN and Bioconductor are
preferable to distribution platforms like GitHub due to this
complication {[}20,37,40{]}. Further, when a package is installed from
GitHub, typically only the source code is available, not binaries
pre-built for your computers operating systems. This means that the
installation process will need to build the package from source on your
computer, which may require certain compilers and other tools (e.g.,
XCode for macOS, Rtools for Windows).

When you're searching for and selecting packages, it's important to note
that a single package is often available in multiple places. Packages
that are on CRAN, Bioconductor, and/or rOpenSci are usually on GitHub.
There's an unofficial, read-only GitHub mirror at CRAN,
\href{https://github.com/cran}{METACRAN}, on which \emph{all} CRAN
packages have their own GitHub repositories with all versions, original
dates, and authors. Additionally, package developers frequently use
their personal GitHub accounts to publicize a package while it's being
developed and continue to host development versions after it's been
published on formal repositories such as CRAN. In this case, CRAN would
host a ``stable'' version of the package, which may simultaneously be in
the ``master'' branch of the developer's GitHub repository for the
package; forthcoming package updates would remain in other branches
(e.g., a ``development'' branch) in the same GitHub repository. Although
less relevant to new R users, this means that you can access and test
upcoming versions of a package through GitHub, which may be particularly
helpful if you're developing a package with dependencies and want to
know if upcoming changes to a package that your code depends on will
cause issues. Lastly, there are some instances where a packages is
posted on CRAN---not only because CRAN is its main repository, but also
to ease its connection with base R---then also published in a
specialized repository. As part of their review process, for example,
rOpenSci asks authors of submitted packages if they plan to publish on
CRAN (or Bioconductor), and packages often become available in both
places.

\hypertarget{rule-4-explore-the-availability-and-quality-of-help}{%
\section{Rule 4: Explore the availability and quality of
help}\label{rule-4-explore-the-availability-and-quality-of-help}}

Current sources of information related to R packages are dispersed and
plentiful. On one hand, this allows users to explore diverse solutions
and discover new tools; on the other, not knowing where to find help can
lead to inefficient and ineffectual roundabouts. Not all package
resources share the same level of quality and the fact that there are
many resources in aggregate does not imply that every package is
associated with the same availability of resources.

Though all published R packages meet some minimal standard of
documentation, beginners and users of complex packages will often want
help documentation that is more thorough than this minimum. Help
associated with R packages generally falls into one of two categories:
help documentation that comes embedded in the source code of the package
(i.e., ships with the package) and documentation or other resources that
are available more broadly, beyond the package's source code. The
documentation that ships with the package can consist of help files,
vignettes, \texttt{README} files, tutorials, and \texttt{pkgdown}
documentation {[}41{]}. Broader help resources can include webpages or
online books whose content is hosted outside the package source, online
tutorials, electronic mailing lists, and large and vocal communities of
R users. Implicitly, if an author has expended time and effort to
provide copious and useful help documentation---within package code,
outside of it, or both---beyond the utility of the resources themselves,
this suggests that the authors are serious about their package
development and adds confidence that they know what they (and their
package) are doing. Exemplary documentation can signify an exceptional
package, while you might keep looking if the quality and quantity of
help resources for a package is lacking.

At the most granular level, every R package should come with a help file
for each outward-facing function---that is, functions which the package
maintainer intends for you to use directly (as opposed to smaller help
functions meant solely for use within other package functions). This
practice is required for packages on some repositories, like CRAN, but
might not be adopted for those posted on less restrictive platforms,
like GitHub. For each outward-facing function in a package, once you've
loaded the package in your R session, you can access its help file from
the console via the \texttt{help} function or by typing \texttt{?}
followed by the function's name (e.g., \texttt{help(mean)} or
\texttt{?mean} returns the help file for the \texttt{mean} function).
Function help files provide information about how to use the function,
including the parameters that can be set, the format to use when
specifying arguments for these parameters, a description of the file,
and even references to papers or other resources on which the function's
underlying algorithm was based. They often contain executable examples
to help you get a feel for how to use the function in practice. If you'd
like to check all help documentation for a package, you can access an
index of help pages from the console via
\texttt{help(package\ =\ "...")} {[}9{]}. The documentation that
accompanies functions within packages is critical; the fact that anyone
can read it anytime and use it to guide their own work facilitates
extensibility. The extent and quality of content in these help files
will vary; ideally, packages should have thorough documentation at the
function level. If the help pages alone leave you wanting, it may be
less likely that the package has further (quality) documentation and
therefore you might reconsider whether that package should be your first
choice, if comparable options exist. In short, if the developer cannot
initially communicate how their tool works, then you probably won't want
to use it in your kitchen (read: if the instruction manual is useless,
don't use the blender).

Fortunately, plenty of R packages include additional documentation
beyond this minimum. While the function-level help files can be rather
technical and extraneous to new users, a vignette is a practical type of
documentation in the form of a tutorial. A vignette is a detailed,
long-form document that describes the problems an R package can solve,
then illustrates applications through clear examples of code with
coordination of functions and explanations of outcomes. See, for
example, the extensive vignette for the
\href{https://cran.r-project.org/web/packages/sf/vignettes/sf1.html}{\texttt{sf}
package}{[}42{]}. Packages can have multiple vignettes; you can view or
edit a specific vignette or obtain a list of all vignettes for a package
via the \texttt{vignette()} function from base R {[}9{]}. For a package
shared through CRAN or Bioconductor, a list of the package's vignettes
can also be found on the respective homepage for the package.
Particularly for packages hosted on GitHub, developers often create a
mini-vignette in the \texttt{README} to provide a high-level overview of
their package and its components, along with a basic introduction and
short tutorial for its usage. This mini-vignette, displayed on the
package repository's main page, should simply tell you why you should
use it, how to use it, and how to install it {[}24{]}.

Interactive tutorials, most commonly those powered by the
\href{https://rstudio.github.io/learnr/}{\texttt{learnr} package}, are
another effective and popular way to facilitate learning about R
packages {[}43{]}. They're either conceptual and exploratory or specific
and structured in their content and can include a range of features:
narratives, figures, videos, illustrations, equations, code exercises,
quizzes, and interactive Shiny elements. You can run
\texttt{learnr::run\_tutorial(name\ =\ "\ldots{}",\ package\ =\ "\ldots{}",\ shiny\_args\ =\ "\ldots{}")}
to deploy the \texttt{learnr} tutorial included within a package, if one
exists {[}43{]}. Alternatively, as of RStudio v1.3, you can access
\texttt{learnr} tutorials for packages you've installed, directly from
the Tutorial pane in the IDE {[}44{]}. Some developers publish their
\texttt{learnr} tutorials as Shiny apps, which can be called from within
the \texttt{learnr::run\_tutorial} function and accessed outside of the
RStudio interface {[}43{]}.

Several thousand developers have been using \texttt{pkgdown} to create
static HTML documentation {[}41,45{]}. This tool allows developers to
collect dispersed forms of documentation for their package such as help
files and vignettes, then build a single package website with this
information plus more. The
\href{https://pkgdown.r-lib.org/}{\texttt{pkgdown}} website itself as
well as those for
\href{http://naniar.njtierney.com/index.html}{\texttt{naniar}},
\href{https://rnabioco.github.io/valr/}{\texttt{valr}}, and
\href{https://xavierdidelot.github.io/TransPhylo/articles/TransPhylo.html}{\texttt{TransPhylo}}
are examples of \texttt{pkgdown} documentation in action for R packages
{[}41,46--48{]}. Sometimes, you can find a link to a package's
corresponding \texttt{pkgdown} website in the ``URL'' section of its
CRAN page, which is often listed alongside its GitHub link.

So far, we've only mentioned types of help documentation that are
incorporated in the source code of a package; they ``ship'' with the
package and so you can investigate them by looking through the
underlying package code. However, there are easier avenues for accessing
help documentation encoded in a package. RDocumentation, for example, is
a searchable \href{https://www.rdocumentation.org/}{website}, package
(\texttt{install.packages("RDocumentation")}), and
\href{https://www.rdocumentation.org/docs/}{JSON API} for obtaining
integrated documentation for packages that are on CRAN, Bioconductor,
and GitHub {[}9,49{]}. RDocumentation may include: an overview,
installation instructions, examples of usage, functions, guides, and
vignettes. This is a subsidiary reason why is may be useful to look for
packages shared on these platforms (see Rule 3).

Some packages have a comprehensive set of resources beyond those encoded
in the package source code. For instance, as mentioned in Rule 2, there
are online and/or print books associated with certain packages. One
popular method that developers use to publish books about their package
is through \texttt{bookdown}, an extension of R Markdown that's
structured to integrate code, text, links, graphics, videos, and other
content in a format that can be published as a free, open, interactive,
and downloadable online book {[}50{]}. The \texttt{bookdown} package
itself has an \href{https://bookdown.org/yihui/bookdown/}{online book}
that details usage of the package {[}50,51{]}.

Many other forms of help documentation can be shared outside of the
package's code. Some packages have corresponding online courses, which
might be taught by the maintainer.
\href{https://noamross.github.io/gams-in-r-course/}{\emph{GAMs in R}},
for instance, is a free interactive course on Generalized Additive
Models in R with the \texttt{mgcv} package {[}52{]}. Every now and then,
an online course will cover one or a few select packages that revolve
around a central theme such as a types of modeling. A suitable example,
\href{https://supervised-ml-course.netlify.app/}{\emph{Supervised
Machine Learning Case Studies in R}}, features several \texttt{tidy}
tools with a primary focus on \texttt{tidymodels} {[}53{]}. Galleries,
as the name would imply, are sites that display works, articles, and/or
code examples associated with certain packages. As you might expect,
galleries can be valuable sources of inspiration, particularly for
packages that provide tools for data visualization tasks. See, for
example, \href{https://www.r-graph-gallery.com/}{The R Graph Gallery} as
well as the galleries for
\href{https://exts.ggplot2.tidyverse.org/gallery/}{\texttt{ggplot2}
extensions},
\href{https://rstudio.github.io/dygraphs/index.html}{\texttt{dygraphs}
for R}, \href{https://shiny.rstudio.com/gallery/}{\texttt{shiny}}, and
\href{https://gallery.htmlwidgets.org/}{\texttt{htmlwidgets}} for
organized displays of outputs from these packages {[}26,54,55{]}. The
\href{https://gallery.rcpp.org/}{\texttt{Rcpp} gallery} is an instance
of a gallery that contains collections of articles and code examples for
the \texttt{Rcpp} package {[}56{]}. Cheatsheets provide concise usage
information for popular packages through code and graphics organized by
purpose. RStudio hosts a number of these, including cheatsheets made by
RStudio's team and those contributed by others. These cheatsheets can be
accessed directly via the RStudio Menu (Help \textgreater{} Cheatsheets)
or from the \href{https://rstudio.com/resources/cheatsheets/}{RStudio
website} on which you can subscribe to cheatsheet updates and find
translated versions. Packages with their own cheatsheets include:
\texttt{caret} (statisical modeling and machine learning),
\texttt{data.table} (data manipulation), \texttt{devtools} (package
development), \texttt{dplyr} (data transformation), \texttt{ggplot2}
(data visualization), \texttt{leaflet} (interactive maps),
\texttt{randomizr} (random assignment and sampling), \texttt{sf} (tools
for spatial data), \texttt{stringr} (character string manipulation), and
\texttt{survminer} (survival plots) {[}11,26,39,42,57--62{]}. Further
information about some R packages is available in video tutorials,
webinars, and code demonstrations (i.e., ``demos'').

There are also resources that involve \emph{asking} for help, and you
can size up the presence of the package you're considering through these
resources, as well. The \emph{R-help} mailing list, to which you can
subscribe and send questions, is moderated by the R Core Development
Team and includes additional facets for major announcements about the
development of R and availability of new code (\emph{R-announce}) and
new or enhanced contributed packages (\emph{R-packages}) {[}63{]}. In
the early days, \emph{R-help} was the only way to seek direct assistance
outside of knowing and asking other R users personally; since, the R
community (e.g., \href{https://community.rstudio.com/}{RStudio
Community}, \href{https://support.bioconductor.org/}{Bioconductor
Support}) has evolved to make asking for help more widespread and
accessible, with inclusion and creative problem-solving as hallmarks of
its online presence {[}64{]}. Certain packages have independent mailing
lists; \texttt{statnet} is an example of a suite of packages that has
its own \href{http://statnet.org/}{community mailing list} {[}65{]}. If
a package has a development repository on GitHub, check the Issues to
verify that the maintainer is responsive to posts and fixes bugs in a
timely manner. In addition, you can search discussion forums such as
\href{https://stackoverflow.com/}{Stack Overflow},
\href{https://stats.stackexchange.com/}{Cross Validated}, and
\href{http://www.talkstats.com/}{Talk Stats} to assess the activity
associated with the package in question. Analyses of the popularity of
comparable data analysis software in email and discussion traffic
suggest that R is rapidly becoming more prevalent and is the leading
language by these metrics {[}6,66{]}. When you encounter a problem, it's
good practice to first update the package to see if the trouble is due
to a bug in a previous version---if the problem persists, seek help by
finding or posting a reproducible example {[}67{]}.

As part of learning to assess the help resources for any package, it's
helpful to reference a very well-documented package. One example of a
package---mentioned a few times already---with notable documentation and
first-rate help resources is \texttt{Rcpp} {[}56{]}. The developers
maintain both a \href{http://www.rcpp.org/}{main} and
\href{http://dirk.eddelbuettel.com/code/rcpp.html}{additional} website
with a wealth of organized information about the package and resources,
including a gallery of examples, associations, publications, articles,
blogs, code, books, talks, and links to other resources with
\texttt{Rcpp} tags. The package has a dedicated
\href{https://lists.r-forge.r-project.org/cgi-bin/mailman/listinfo/rcpp-devel}{mailing
list}, for developers using the package. \texttt{Rcpp} is also
thoroughly described in a print book by its maintainer {[}68{]} and in
chapters or sections of other books (e.g., {[}67{]}).

\hypertarget{rule-5-quantify-how-established-it-is}{%
\section{Rule 5: Quantify how established it
is}\label{rule-5-quantify-how-established-it-is}}

Consulting data to inform comparisons is never a bad idea; numerical
data associated with R packages will give you an impression of how
popular a tool is and whether it's stood the test of time. There are
metrics that you can use to determine how often an R package has been
downloaded, how long it's been around, and how often (and recently) it's
updated through new versions.

You can find several of these metrics by visiting the site(s) on which a
package is shared (see Rule 3). If a package is shared through a GitHub
repository, the homepage for the repository will list the number of
GitHub Stars, Forks, and Watchers for that particular repository. As
proxies for a repository's overall following---and possibly indirect
measures of quality and impact---Stars indicate that people like or are
interested in a project, Forks reveal the number times a repository has
been copied, and Watchers represent GitHub users who are subscribed to
updates on a repository's activity {[}69{]}. Although not a definitively
unbiased metric, a large number of Stars, Forks, and Watchers associated
with a package implies a substantial following and widespread usage
{[}70{]}. If a package is posted on CRAN or Bioconductor, its page on
that site will include the date on which its latest version was released
as well as a link to ``Old sources''---previous versions of the package
and the release date of each. The page also provides a link to a
\texttt{NEWS} file, where you can see a list of all published versions
as well as the major changes made from one version to the next. To
further determine if a package is well-established in the R community,
refer to the number of versions and updates (more is better) in addition
to the date of the most recent versions and updates (newer is typically
better, as an indication of active maintenance, although in some cases a
package with a stable algorithm and few dependencies may not need many
updates). You can also see, on a package's CRAN or GitHub page, the
number of reverse dependencies it has. When a package has numerous
reverse dependencies, the code underlying many other R packages
incorporates functions from that package, which can indicate that other
R developers find that package useful.

RDocumentation (see Rule 4) is rich with stats on R packages, uniting
many of the previously discussed metrics, while also providing
cross-package summaries. RDocumentation hosts a live
\href{https://www.rdocumentation.org/trends}{Leaderboard} with trends
including the number of indexed packages and indexed functions, most
downloaded packages, most active maintainers, newest packages, and
newest updates. What's more, each package is assigned a percentile
rank---featured on its RDocumentation page---that quantifies the number
of times a package has been downloaded in a given month. A ranking
algorithm computes the direct, user-requested monthly downloads by
accounting for reverse dependencies (indirect downloads) so packages
that are commonly used within other packages, and hence frequently
downloaded as automatic downloads with those packages, don't skew the
calculation {[}71{]}. You can research stats on corresponding
dependencies for a more holistic picture.

Most of the above-mentioned resources will help you learn more about
packages on an individual basis. Wouldn't it be nice if you could
compare related packages that may have features that overlap? Some
clever R users at an Unconference thought so, too, and developed a
package called
\href{https://docs.ropensci.org/packagemetrics/}{\texttt{packagemetrics}:
A Package for Helping You Choose Which Package to Use}, hosted at
rOpenSci {[}72,73{]}. You can install \texttt{packagemetrics} via
\texttt{devtools::install\_github("ropenscilabs/packagemetrics")} and
run \texttt{packagemetrics::package\_list\_metrics()} where the input is
a vector of packages of interest, by name {[}39,72{]}. Then, you can use
\texttt{packagemetrics::metrics\_table()} to obtain a table of
information that includes each package's respective number of downloads,
GitHub Star rating, \texttt{tidyverse} compatibility, and several more
indicators to help you decide which package is most established and
widely used (Fig \ref{fig7}) {[}72,74{]}.

\begin{figure}[!h]
\caption{{\bf Output from the \texttt{packagemetrics} package, used here to compare data visualization packages}
The R package, \texttt{packagemetrics}, is a convenient tool for directly comparing available packages for a task of interest, using metrics related to package development, documentation, and compatibility with other packages.}
\label{fig7}
\end{figure}

An important caveat to consider: several of the metrics introduced in
this Rule are based on how often a package is downloaded; however, it's
not uncommon for users to download a package to explore its
functionality and ultimately decide against using it for their task.
Therefore, these suggested metrics can provide a first pass at
identifying popular packages---and will, in particular, help you
distinguish popular packages for common general tasks, like data
wrangling and visualization---but may be less helpful when you're
selecting a package for a very specialized task. In this case, evidence
of peer acceptance and review will usually be more instructive, which
brings us to our next Rule.

\hypertarget{rule-6-seek-evidence-of-peer-acceptance-and-review}{%
\section{Rule 6: Seek evidence of peer acceptance and
review}\label{rule-6-seek-evidence-of-peer-acceptance-and-review}}

Peer review is an important aspect of scientific research, not least
because it establishes scholarly credibility. While scientific articles
undergo a peer review process that's fairly standardized across
disciplines, R packages have a broader range of possibilities for peer
review or other evidence of acceptance by members of the scientific
community. You can research information about an R package in different
forms of literature and determine the extent to which it's been
validated by the scientific community.

First, some maintainers of select packages have written an article
describing their package and its algorithms; such articles must go
through the traditional peer-review process in order to be published.
Certain journals publish articles about R packages themselves while
others feature work that used a particular package. Journals with a
targeted and technical focus on R packages and other open source
scientific software include the
\href{https://joss.theoj.org/}{\emph{Journal of Open Source Software}},
the \href{https://www.jstatsoft.org/index}{\emph{Journal of Statistical
Software}}, and \href{https://journal.r-project.org/}{\emph{The R
Journal}}. Discipline-specific journals like
\href{https://bmcbioinformatics.biomedcentral.com/}{\emph{BMC
Bioinformatics}}, may also publish articles to introduce R packages that
have proven to be useful for aspects of the bioinformatics research
process, for example. Try searching various websites for prominent
journals in your field with queries that include \texttt{"R\ package"}
to find these types of articles. The number of Google Scholar citations
for a package or, similarly, an article about a package, can also help
serve as a metric for a package's impact on scientific research and
utility in research contexts (see Rule 5).

For an R package featured in a peer-reviewed journal article, you can be
assured that the ideas and principles behind the package have been
evaluated by outside reviewers. However, such a review may not have
included a fine-grained peer review of the underlying code in the
package itself. There's an expanding movement toward direct peer review
of R packages and other open source scientific software at the code
level. More specifically, rOpenSci is a unique example of an ecosystem
of open source tools with peer reviewed R packages (see Rule 3)
{[}31{]}. The rOpenSci organization oversees peer review of R packages
at a highly detailed level that's increasingly desired by scientists;
packages published through their repository have passed external review
of the code. They limit their scope to packages that help scientists
implement reproducible research practices and manage the data lifecycle
{[}32{]}. Their review process is conducted openly on GitHub and, as a
result, the scientific community can view the discussions between
package maintainers and peer reviewers that form this transparent
process.

You can also seek evidence of a package's peer acceptance by determining
how often researchers have used it within scientific papers describing
research, rather than package-focused papers. Useful packages may be
mentioned in the Methods and/or References sections of such papers.
Packages hosted on CRAN or Bioconductor have a standardized format for
citation, and a call to \texttt{citation("package")} will print the
suggested citation for the package {[}9{]}. Further, in response to the
rising number of researchers creating tools and software to work with
their data, GitHub has granted developers the ability to obtain a
Digital Object Identifier (DOI) for any GitHub repository archive so
that code can be cited in academic literature {[}75{]}. You can search
for packages directly by name in Google Scholar: the \texttt{Cited\ by}
link displays the number of times a package has been cited and connects
to a page with those publications base on citations using these formats.
Further, if a package has such a DOI, you can explore the network of
research associated with that package. While it's best practice to cite
R packages using these citation formats, occasionally R packages will be
mentioned in the Methods section of a paper without a formal citation.

Finally, many R packages are associated with content in books and series
from scientific publishers. For example, Springer has a
\href{https://www.springer.com/series/6991}{\emph{Use R!} series} with
books covering specialized topics from finance to marketing, from
chemometrics to brain imaging, and more. Additionally, consider reading
titles from
\href{https://www.routledge.com/Chapman--HallCRC-The-R-Series/book-series/CRCTHERSER}{\emph{The
R Series} at Chapman \& Hall/CRC}. The books in this series span topics
related to R in terms of its applications (e.g., epidemiology,
engineering, social sciences), uses (e.g., statistical methods), and
development (e.g., creating packages).

\hypertarget{rule-7-find-out-who-developed-it}{%
\section{Rule 7: Find out who developed
it}\label{rule-7-find-out-who-developed-it}}

There's variation in R users' backgrounds and skills, and the same is
true for R developers. That said, the R community prides itself on
embracing newcomers at all levels of involvement. This Rule doesn't
imply that worthwhile packages are exclusively written by well-known
authors. Rather, associations and reputation can be a proxy for quality;
in this way, the process of evaluating and comparing R packages is no
different than other decisions. In fact, as you become more immersed in
the R community, you'll find that name recognition is a factor, among
many, that helps you establish trust in certain tools more quickly than
others {[}70{]}.

Information about the package authors can be found in the metadata of
the package. If you're looking at the source code of the package, there
should always be a \texttt{DESCRIPTION} file with metadata. The
\texttt{DESCRIPTION} is a succinct record of the package's purpose,
dependencies, version, date, license, associations, authors, and other
technical details. In the ``Author'' field of this file, you can
identify both the ``Maintainer'' of the package (a similar role to that
of the corresponding author for an article) as well as other authors or
contributors. For packages hosted solely on a platform like GitHub, this
might be the only way to discover a package's main author, but profile
links for all authors are available under the Contributors section on
the package's main GitHub page. On the other hand, for packages hosted
on repositories, this package data is usually also posted by the
repository somewhere online. Packages on CRAN have a page of metadata
that's accessible through a standard web address (see Rule 3). Here, you
might notice ORCiD links and/or various abbreviations associated with
the names in the ``Author'' and Maintainer" fields that denote roles
such as author \texttt{{[}aut{]}}, creator/maintainer
\texttt{{[}cre{]}}, contributor \texttt{{[}ctb{]}}, and funder
\texttt{{[}fnd{]}}. Similarly, Bioconductor packages each have their own
webpage, with some metadata at the top of the page, including the
maintainer and other authors along with their roles (see Rule 3).

You can assess the credibility and commitment of R package developers
through direct and indirect signals. Who made the package? Consider
whether expertise in a certain domain is vital to the design and
creation of the tool. Research the authors' associations in academia,
government, and/or industry and gauge the extent to which they have a
primary role in R development. Sometimes, it's possible to deduce author
affiliations by simply looking at their domain name on their email
addresses. Affiliations may serve as a proxy for package quality and
sophistication insofar as standards, on average, tend to be higher in
regulated settings. R packages produced by the
\href{https://owi.usgs.gov/R/}{United States Geological Survey (USGS)},
for instance, are maintained by domain experts that know field-specific
data sources and algorithms very well. Moreover, authors with
affiliations that include R package development as a paid responsibility
(e.g., USGS, RStudio) might, on average, produce better packages. If a
package was developed by multiple authors, research each of them to
evaluate the robustness of the team. You can learn more about their
experience, active contributions to the R community, and history related
to package development by exploring their profiles on GitHub, Google
Scholar, Research Gate, Twitter, or personal or package websites.
Frequent commits and effective resolutions of GitHub Issues can be a
testament to the authors' priorities and commitment. If an author has a
history of package development, peruse their portfolio of packages to
see if any are highly regarded, recognizable, or impressive at first
glance. When you use a great tool in your kitchen, you consider buying
more tools from that brand. If, in the past, you found a package that
you like to use and is helpful for your work, it can be worthwhile to
research other packages made by that author and/or maintainer. You might
benefit from making a habit of paying attention to who makes tools that
work well for your needs and are designed in a way that seems intuitive
to you---take note of these ``favorite'' package authors.

As you may know, the R community values collaboration---it's not
uncommon for packages to be created and maintained by more than one
person and, in many cases, large teams. Roles within these teams can
shift as the package evolves. As an example, \texttt{ggplot2} has no
small team of authors and its original maintainer isn't the same as its
current one {[}26{]}. Oftentimes, existence of a large team can give you
clues (although not guaranteed) about how the package is evolving in the
R ecosystem. For one, it could indicate that a package is very useful
and has generated rich community involvement, as is the case for
\texttt{ggplot2}---many of its authors were added after the first
release {[}26{]}. This might reveal that a package is part of a
larger-scale and more organized project. Also, a sizeable team could
imply that the package grew out of a substantial collaboration, like a
funded grant with a multi-person team, which happens a lot for specific
methods packages. The \texttt{flowCore} package is one example of such a
case {[}76,77{]}. This might suggest that a package has a certain level
of review and investment from several R programmers (the authors), not
just one, as in a solo-authored project. It's also useful to examine the
full list of authors (even if it's long) because you might not recognize
the maintainer's name, but you might recognize another author's name.
This can happen if the maintainer is a student or someone else that's
up-and-coming; sometimes, more senior authors don't assume the role of
maintainer, even though they're quite invested in the project. It's
worth looking to see, for example, if someone who is well-known for
excellence in methods development is one of the non-maintaining authors,
which may help you gauge the quality of the fundamental algorithms and
techniques in the package code.

\hypertarget{rule-8-see-how-its-developed}{%
\section{Rule 8: See how it's
developed}\label{rule-8-see-how-its-developed}}

You don't need to be a software engineer to identify strong package
development. Scientific software developers sometimes neglect best
practices; indeed, these shortcomings are evident in the tools they
create {[}78{]}. There are concrete ways to measure a tool's robustness
beyond whether it works for those who didn't create it. As a user, there
are four main relevant package development protocols---in addition to
some related indicators---that you can investigate to assess the
underlying stability and utility of a package: dependencies, unit
testing, version control, and continuous integration.

R packages often depend on other R packages; you should check the
reputations of such dependencies when selecting a package---quality
packages will rely on a solid web of quality packages. What's more, like
other types of software, well-maintained R packages have multiple
versions corresponding to iterative releases to indicate that the
package is compatible with dependencies and loyally updated (e.g., bug
fixes, general improvements, new functionality) {[}24,79{]}. You can
explore the version history of a package to see if it's up-to-date (See
Rule 5). Package dependencies are not only important to establish
stability, but might also give you a hint as to whether a package will
fit well into your workflow. This is an especially critical
consideration if you're using a ``tidy'' workflow or working with
special object types in, for example, a Bioconductor workflow. You can
review the list of dependencies for a package in question to check if it
depends on some or many of the packages in your workflow. If this is the
case, you can be fairly confident that the package uses the same (or at
least similar) style and conventions as those you want to use.

A responsible developer with a consistent and reproducible workflow will
implement formal testing on their code to examine expected behavior via
an automated process called unit testing {[}24,80{]}. This helps assure
the developer (and you) that the package's underlying units of code do
what they intend to do. Although inconvenient at the outset, the
developer---and by extension, the package user---will benefit from unit
testing, which results in fewer bugs, a well-designed code structure, an
efficient workflow, and robust code that's not sensitive to major
changes in the future {[}24{]}. To alleviate the burdens of unit
testing, \texttt{testthat} is a popular, integrative R package that
helps developers create reliable functions, minimize error, and
visualize progress through automatic code testing {[}81{]}. Developers
are also interested in quantifying the amount of code in their package
that has been tested. Test coverage, a measurement of the proportion of
code that has undergone unit testing, is an objective metric for package
developers, contributors, and users to evaluate code quality. Many
developers use the \texttt{covr} package to generate reports and
determine the magnitude of coverage on the function, script, and package
levels {[}82{]}. You can look at the source code for a package to
confirm it has a testing suite; for instance, if you're on a package's
GitHub page, you can navigate to the main Code page and look for a
folder labeled \texttt{tests} or something similar.

Relatedly, developers who host their packages on GitHub often post
status badges in the overview (\texttt{README}) section of the
repository webpage for the master branch. GitHub badges are a common
self-imposed method to signal use of best practices and motivate
developers to produce software that's high in quality and transparency
{[}83{]}. They're easily visible, interactive, and increase readability
of the \texttt{README}. You might see, for example, license
(\texttt{license}), version (\texttt{release}), dependency
(\texttt{dependencies}), code coverage (\texttt{codecov}), or build
status (\texttt{build}) badges, all of which are good indicators of
package caliber. What's more, badges have dynamic status tags to
indicate the current status of the particular part of a workflow to
which the badge refers; as an example, a \texttt{build} badge with a
\texttt{passing} status tag could be telling you that the package is
passing continuous integration (CI) tests (CI is explained later in this
Rule) {[}69{]}. Badges are also clickable buttons that link to a page
that describes and displays detailed aspects of the metric associated
with the badge. For instance, when you click on a \texttt{codecov}
badge, you'll likely end up on a page with a timeline, graphics, and
information about the percentage and distribution of test coverage for
the package. You can refer to special badges like \texttt{CRAN},
\texttt{Bioconductor}, or \texttt{rOpenSci} to determine which
repository (or repositories) a package is shared on (See Rule 3).

As we mentioned in Rule 3, version control has an essential role in
package development and computational literacy more broadly {[}24,84{]}.
Version control is like a time capsule for your workflow because it
monitors and tracks changes to files as a project evolves, and stores
them as previous versions to be recovered if necessary. In other words,
``version control is as fundamental to programming as accurate notes
about lab procedures are to experimental science'' {[}{[}84{]}; p6{]}.
Git is a decentralized open source version control system that's useful
regardless of whether a project is independent or collaborative
{[}85{]}. GitHub works in conjunction with Git to provide a powerful
structured system to organize and manage components of a project for
others and your future self. A growing number of scientists have
research programs based in GitHub, which has become a revolutionary tool
for productive team science and distributed development efforts
{[}79,86{]}. As you may expect, Git coupled with GitHub is the version
control duo of choice among serious R package developers {[}24{]}. Thus,
if the package you're interested in using is among the thousands hosted
on GitHub, this is evidence that the developer is at least committed to
a logical, open, and reproducible workflow, suggestive of more time
spent designing their tool.

Lastly, developers often implement CI in conjunction with using a
version control hosting platform like GitHub. For R package developers,
CI typically involves utilizing an outside service, like
\href{https://travis-ci.org/}{Travis CI}, to run automated tests (e.g.,
unit tests, CRAN checks) on package code every time they push new
contributions and changes to their (potentially shared) GitHub
repository {[}24,32{]}. Just as you're happy to know that a manufacturer
guarantees their kitchen tool is a unified realization of many working
parts, you should be impressed when a package is set up to use CI. This
is particularly important for packages with a multi-person development
team, but also an indicator of quality for single-author packages. Use
of CI is a strong sign that the components of the package are regularly
checked to ensure that code passes tests and the package builds without
failures {[}87{]}.

\hypertarget{rule-9-put-it-to-the-test}{%
\section{Rule 9: Put it to the test}\label{rule-9-put-it-to-the-test}}

If you're unable to decide whether to use a package based on prior
Rules, test it out. Similarly, if you have narrowed your options, work
with each to highlight differences. Exploring the package and engaging
in trial and error using your skills in context of your goal will
illuminate technical details and solidify any doubts.

At this point, what you've learned about the package should be quite
helpful. If the package has a vignette, this is a great place to start
experimenting with code. Vignettes include many common data science
problems with solutions; you can run the code examples, tweak them, and
compare the outputs. First, try some examples with the example data
provided by the package, then see if you can translate that to working
with your own data. Sometimes, package functions will require inputs to
be in a specific data type (e.g., a dataframe, or a certain type of S4
object) or to have a specific format (e.g., if it's in a dataframe, to
have certain columns with certain names). By moving from the example
data to your own, you can gauge the effort it will take to transform
your data to fit these formats. You can also be on the lookout for
whether certain characteristics of your data (e.g., fields with missing
values) cause complications when using the package's functions. You can
continue by looking through the collection of help files for the
package. Help files often include executable examples, and you can again
start by testing these with the example data provided by the package,
then move on to trying them with your own data and within your own
workflow.

Sometimes packages don't interact well with other packages; a recipe
prepared with an odd combination of tools won't turn out. You should
therefore test the \emph{interoperability} of all the packages you want
to use. There are some packages that are masterful at doing what they're
made to do, yet incongruous with other packages. Such packages might,
for example, use S3 or S4 objects, which are two main approaches
developers use to implement object-oriented programming in R. Many
packages for spatial analysis, as well as those from Bioconductor, tend
to use S4 objects to represent data {[}7{]}. On the other hand, the
\texttt{tidyverse}, a unified suite of packages with an grammatical
structure employed within a ``pipeline,'' expects dataframe objects
{[}74{]}. Because of this, when you're working in the
\texttt{tidyverse}, you may need to take extra steps to extract
components from S3/S4 objects to re-frame the data into a dataframe.
Some packages already exist to help with this interfacing task between
S3/S4 objects and \texttt{tidyverse} tools. The \texttt{broom} package,
and the bioinformatics analog, \texttt{biobroom}, aim to alleviate these
disruptions by converting ``untidy'' objects in S3/S4 objects into
``tidy'' data in dataframes, thereby making it easier to integrate
statistical functions into the structure of the \texttt{tidyverse}
workflow {[}74,88,89{]}. Similarly, the \texttt{caret} package
facilitates interoperability for machine learning packages by providing
a uniform interface for modeling with various algorithms from different
packages that would otherwise have independent syntax {[}57{]}.

\hypertarget{rule-10-develop-your-own-package}{%
\section{Rule 10: Develop your own
package}\label{rule-10-develop-your-own-package}}

What if you can't find a package to solve your data science problem?
Though you may not be the best cook, we all come upon times---of
creativity or necessity---that prompt us to experiment in our kitchens.
Rather than exclusively being a user of packages, you can create
them---more easily than you may think {[}24{]}.

If you decide to write your own R code to perform your task (see Rule
1), a natural next step is to wrap this code into a function that can be
re-used; an R package is basically a collection of one or more of these
functions. The reasons why you might want to create a package are
abundant: necessity, innovation, standardization, automation, specialty,
containment, organization, sharing, collaboration, extensibility, and
more. Whatever your motivation, packages are simply sets of tools; you
can create a package out of any collection of specialty functions. A key
tenant of efficient programming is the ``DRY'' principle---Don't Repeat
Yourself. Functions help with this and are warranted when you copy and
paste code while making slight modifications after each iteration
{[}67{]}. You should use a function when you want to replace repetitive
code with a defined, compact piece of code defined only once in your
script, which you can call throughout in shorthand. Since R packages are
just sets of these functions, they're not solely reserved for
large-scale, comprehensive methods; rather, package development can also
help you learn how to apply proper coding techniques to writing
functions and documentation with reproducibility and collaboration in
mind {[}90{]}. Packages need not be formal nor entirely cohesive. For
instance, personal R packages such as \texttt{Hmisc} and
\texttt{broman}, are comprised of miscellaneous functions which the
creator has developed and frequently uses {[}91,92{]}. They're like the
collection of little tried-and-true tricks and tools you've made to cook
your own recipes in your kitchen.

While personal R packages are virtually unbounded, we'll describe some
of their immediate uses. If you often do the same specific,
discipline-oriented data analysis tasks---for which either a package
doesn't exist or you don't want to rely one that does (See Rule 1)---a
personal R package could help you automate some redundant aspects of
your work. You can even connect components of other tools you're using,
that wouldn't otherwise play nicely together. For example, if you're
work largely involves packages that output S3/S4 object types, but you
want to use standard \texttt{ggplot2} tools to visualize the output,
you'll encounter a problem because functions in \texttt{ggplot2} expect
a dataframe as the standard input, not these more complex object types
{[}26{]}. In this case, you might want to write some of your own
functions that extract what you need from the S3/S4 objects and convert
them into a dataframe input type so they're compatible with
\texttt{ggplot2}. Interestingly, this type of occurrence is the precise
motivation behind the widely used \texttt{broom} package {[}88,93{]}.
You can also create visualization tools that are tailored to your
research and include them in your personal R package. If you frequently
design a certain type of plot with your data, you can encode that
process into a function, even if it relies heavily on functions from
other packages. For instance, this could happen if you implement
particular tweaks whenever you make scatterplots with your data; many R
users build custom themes within \texttt{ggplot2} for purposes like
personal or professional branding, which can be wrapped into functions
and contained in a personal R package {[}26,87{]}.

Although you may not anticipate that anyone else will use your tools,
following best practices for package development is still a good idea.
You wouldn't want to craft a useful cooking hack without having a good
way to replicate it when you need to cook the same (or similar) recipe
again. Even if you're the only person who will use your package, be kind
to your future self---as a consumer of shared packages, you know the
inherent benefits of robust software development relative to the quality
of code, data, documentation, versions, and tests {[}78{]}. There's also
no reason not to take advantage of existing resources when coding your
package; all your newfound special tools paired with others' shared
tools can help you make exquisite recipes. There are R packages that aid
in package development (e.g., \texttt{devtools}, \texttt{usethis},
\texttt{testthat}, \texttt{roxygen2}, \texttt{rlang}, \texttt{drat})
{[}33,39,81,94--96{]}---while these packages can provide
industrial-strength support for developing large and complex packages,
they also prove useful even in developing small, private packages. They
make it easier and more efficient, for example, to write help
documentation and unit tests for functions in your package. RStudio's
IDE has a framework called ``Projects,'' which can help you establish an
organized framework of directories within an R Project
(\texttt{.Rproj})---specifically designed for creating R
packages---that's straightforward to navigate throughout the package
development process. Tips on using many of these package development
tools are available through resources authored by expert R developers,
like \href{https://r-pkgs.org/}{\emph{R Packages}} by Hadley Wickham
{[}24{]}, while the official resource on writing packages to share
through CRAN is the manual,
\href{https://cran.r-project.org/doc/manuals/r-release/R-exts.html}{\emph{Writing
R Extensions}} {[}97{]}.

Regardless of whether you're writing a package to share, or one for your
own use that you don't plan to publish through a repository, consider
using version control. Among the numerous reasons why version control
will positively impact your workflow, you'll be able to easily share
your project with collaborators, at any point in time. If you decide to
collaborate on a package, even if you just intend for the package to
remain a resource among your group, the Project framework is compatible
with distributed development---a feature that couples well with version
control and online platforms for version controlled repositories, like
GitHub. Version control also helps you maintain a clean directory of
code without limiting the process of making changes to files. All
changes are trackable and rewindable so earlier versions are restorable.
You can create offshoots of your main project (branches) to try new
things and incorporate some (or all) of the changes into the main
branch. As you work, keep in mind that within the open source R
community, there's no lack of effective organizational frameworks to
reference; in fact, repositories for many exemplary packages are readily
available on GitHub.

\hypertarget{conclusion}{%
\section{Conclusion}\label{conclusion}}

The open source framework is the heartbeat of R; it sustains and
enlivens the growth and improvement of the R ecosystem. Anyone, whether
they're motivated to propel a research program or interested in learning
from data in other contexts, can be an R user---without significant
barriers to entry. A number of R users don't stop there, however. They
become R makers by adding their own contributions so not only they, but
new and existing users and makers, can do more in R. This plays an
essential role in the lifecycle of R's beloved shared tools.

R packages are a defining feature of the language insofar as many are
robust, user-friendly, and richly extend R's core functionality. Some of
the most prominent R packages are a result of the developer abstracting
common elements of a data science problem into a workflow that can be
shared and accompanied by thorough descriptions of the process and
purpose. In this way, R packages have effectively transformed how we
interact with data in the modern day in, perhaps, a more impactful
manner than several revered contributions to theoretical statistics
{[}5{]}. Packages greatly enhance the user experience and enable you to
be more efficient and effective when working with data, regardless of
prior experience.

Nonetheless, the sheer quantity and potential complexity of available R
packages can undermine their collective benefits. Finding and choosing
packages, particularly for beginners, can be daunting and difficult. R
users often struggle to sift through the tools at their disposal and
wonder how to distinguish appropriate usage. These ten simple rules for
navigating the shared code in the R community are intended to serve as a
valuable page in your computing cookbook---one that will evolve into
intuition and yet remain a reliable reference. May searching for and
selecting proper tools no longer spoil your appetite and dissuade you
from discovering, trying, creating, and sharing new recipes.

\hypertarget{acknowledgements}{%
\section{Acknowledgements}\label{acknowledgements}}

The authors wish to thank Bailey Fosdick, Matthew Hughes, Therese
Kondash, and Kathleen Wendt for their helpful input on the manuscript.

\hypertarget{references}{%
\section*{References}\label{references}}
\addcontentsline{toc}{section}{References}

\hypertarget{refs}{}
\begin{CSLReferences}{0}{0}
\leavevmode\hypertarget{ref-rehman2012}{}%
\CSLLeftMargin{1. }
\CSLRightInline{Rehman J. {Open Science and Access to Medical Research}
{[}Internet{]}. Scientific American; 2012. Available:
\url{https://blogs.scientificamerican.com/guest-blog/open-science-and-access-to-medical-research/}}

\leavevmode\hypertarget{ref-peng2011}{}%
\CSLLeftMargin{2. }
\CSLRightInline{Peng RD. {Reproducible Research in Computational
Science}. Science. 2011;334: 1226--1227. }

\leavevmode\hypertarget{ref-goodman2016}{}%
\CSLLeftMargin{3. }
\CSLRightInline{Goodman SN, Fanelli D, Ioannidis JP. {What Does Research
Reproducibility Mean?} Science Translational Medicine. 2016;8:
341ps12--341ps12. }

\leavevmode\hypertarget{ref-stodden2013}{}%
\CSLLeftMargin{4. }
\CSLRightInline{Stodden V, Miguez S. Best practices for computational
science: Software infrastructure and environments for reproducible and
extensible research. Journal of Open Research Software. Ubiquity Press;
2014;2. }

\leavevmode\hypertarget{ref-donoho2017}{}%
\CSLLeftMargin{5. }
\CSLRightInline{Donoho D. {50 Years of Data Science}. Journal of
Computational and Graphical Statistics. 2017;26: 745--766. }

\leavevmode\hypertarget{ref-robinson2017}{}%
\CSLLeftMargin{6. }
\CSLRightInline{Robinson D. {The Impressive Growth of R} {[}Internet{]}.
Stack Overflow; 2017. Available:
\url{https://stackoverflow.blog/2017/10/10/impressive-growth-r/}}

\leavevmode\hypertarget{ref-gentleman2004}{}%
\CSLLeftMargin{7. }
\CSLRightInline{Gentleman RC, Carey VJ, Bates DM, Bolstad B, Dettling M,
Dudoit S, et al. {Bioconductor: Open Software Development for
Computational Biology and Bioinformatics}. Genome Biology. 2004;5: R80.
}

\leavevmode\hypertarget{ref-holmes2018}{}%
\CSLLeftMargin{8. }
\CSLRightInline{Holmes S, Huber W. {Modern Statistics for Modern
Biology} {[}Internet{]}. Cambridge University Press; 2018. Available:
\url{https://web.stanford.edu/class/bios221/book/index.html}}

\leavevmode\hypertarget{ref-Rproject2020}{}%
\CSLLeftMargin{9. }
\CSLRightInline{R Core Team. R: A language and environment for
statistical computing {[}Internet{]}. Vienna, Austria: R Foundation for
Statistical Computing; 2020. Available:
\url{https://www.R-project.org/}}

\leavevmode\hypertarget{ref-hornik2012}{}%
\CSLLeftMargin{10. }
\CSLRightInline{Hornik K. {Are there too many R Packages?} Austrian
Journal of Statistics. 2012;41: 59--66. }

\leavevmode\hypertarget{ref-dplyr}{}%
\CSLLeftMargin{11. }
\CSLRightInline{Wickham H, François R, Henry L, Müller K. {dplyr: A
Grammar of Data Manipulation} {[}Internet{]}. 2020. Available:
\url{https://CRAN.R-project.org/package=dplyr}}

\leavevmode\hypertarget{ref-tidyr}{}%
\CSLLeftMargin{12. }
\CSLRightInline{Wickham H, Henry L. {tidyr: Tidy Messy Data}
{[}Internet{]}. 2020. Available:
\url{https://CRAN.R-project.org/package=tidyr}}

\leavevmode\hypertarget{ref-twitteR}{}%
\CSLLeftMargin{13. }
\CSLRightInline{Gentry J. {twitteR: R Based Twitter Client}
{[}Internet{]}. 2015. Available:
\url{https://CRAN.R-project.org/package=twitteR}}

\leavevmode\hypertarget{ref-mailR}{}%
\CSLLeftMargin{14. }
\CSLRightInline{Premraj R. {mailR: A Utility to Send Emails from R}
{[}Internet{]}. 2015. Available:
\url{https://CRAN.R-project.org/package=mailR}}

\leavevmode\hypertarget{ref-phonenumber}{}%
\CSLLeftMargin{15. }
\CSLRightInline{Myles S. {phonenumber: Convert Letters to Numbers and
Back as on a Telephone Keypad} {[}Internet{]}. 2015. Available:
\url{https://CRAN.R-project.org/package=phonenumber}}

\leavevmode\hypertarget{ref-bioCproject}{}%
\CSLLeftMargin{16. }
\CSLRightInline{Huber W, Carey VJ, Gentleman R, Anders S, Carlson M,
Carvalho BS, et al. {Orchestrating High-Throughput Genomic Analysis with
Bioconductor}. Nature Methods. 2015;12: 115--121. Available:
\url{http://www.nature.com/nmeth/journal/v12/n2/full/nmeth.3252.html}}

\leavevmode\hypertarget{ref-ctv}{}%
\CSLLeftMargin{17. }
\CSLRightInline{Zeileis A. {CRAN Task Views}. R News. 2005;5: 39--40.
Available: \url{https://CRAN.R-project.org/doc/Rnews/}}

\leavevmode\hypertarget{ref-smith2017}{}%
\CSLLeftMargin{18. }
\CSLRightInline{Smith D. {The R Community is One of R's Best Features}
{[}Internet{]}. Revolutions. Microsoft; 2017. Available:
\url{https://blog.revolutionanalytics.com/2017/06/r-community.html}}

\leavevmode\hypertarget{ref-ellis2017}{}%
\CSLLeftMargin{19. }
\CSLRightInline{Ellis SE. {Hey! You There! You are Welcome Here}
{[}Internet{]}. rOpenSci. NumFOCUS; 2017. Available:
\url{https://ropensci.org/blog/2017/06/23/community/}}

\leavevmode\hypertarget{ref-rickert2018}{}%
\CSLLeftMargin{20. }
\CSLRightInline{Rickert J. {What Makes a Great R Package?}
{[}Internet{]}. RStudio; 2018. Available:
\url{https://rstudio.com/resources/rstudioconf-2018/what-makes-a-great-r-package-joseph-rickert/}}

\leavevmode\hypertarget{ref-randomforest}{}%
\CSLLeftMargin{21. }
\CSLRightInline{Liaw A, Wiener M. {Classification and Regression by
randomForest}. R News. 2002;2: 18--22. Available:
\url{https://CRAN.R-project.org/doc/Rnews/}}

\leavevmode\hypertarget{ref-rayshader}{}%
\CSLLeftMargin{22. }
\CSLRightInline{Morgan-Wall T. {rayshader: Create Maps and Visualize
Data in 2D and 3D} {[}Internet{]}. 2020. Available:
\url{https://CRAN.R-project.org/package=rayshader}}

\leavevmode\hypertarget{ref-broman2020}{}%
\CSLLeftMargin{23. }
\CSLRightInline{Broman K. {Getting your R package on CRAN}
{[}Internet{]}. 2020. Available:
\url{https://kbroman.org/pkg_primer/pages/cran.html}}

\leavevmode\hypertarget{ref-wickham2015}{}%
\CSLLeftMargin{24. }
\CSLRightInline{Wickham H. {R Packages: Organize, Test, Document, and
Share your Code}. O'Reilly Media, Inc.; 2015. }

\leavevmode\hypertarget{ref-cranpolicy2020}{}%
\CSLLeftMargin{25. }
\CSLRightInline{{CRAN Repository Policy} {[}Internet{]}. The R
Foundation; 2020. Available:
\url{https://cran.r-project.org/web/packages/policies.html\#Submission}}

\leavevmode\hypertarget{ref-ggplot2}{}%
\CSLLeftMargin{26. }
\CSLRightInline{Wickham H. {ggplot2: Elegant Graphics for Data Analysis}
{[}Internet{]}. Springer-Verlag New York; 2016. Available:
\url{https://ggplot2.tidyverse.org}}

\leavevmode\hypertarget{ref-bioconductor}{}%
\CSLLeftMargin{27. }
\CSLRightInline{{Bioconductor: Open source Software for Bioinformatics}
{[}Internet{]}. Bioconductor; 2020. Available:
\url{https://www.bioconductor.org/}}

\leavevmode\hypertarget{ref-DESeq2}{}%
\CSLLeftMargin{28. }
\CSLRightInline{Love MI, Huber W, Anders S. {Moderated Estimation of
Fold Change and Dispersion for RNA-seq Data with DESeq2}. Genome
Biology. 2014;15: 550.
doi:\href{https://doi.org/10.1186/s13059-014-0550-8}{10.1186/s13059-014-0550-8}}

\leavevmode\hypertarget{ref-BiocManager}{}%
\CSLLeftMargin{29. }
\CSLRightInline{Morgan M. {BiocManager: Access the Bioconductor Project
Package Repository} {[}Internet{]}. 2019. Available:
\url{https://CRAN.R-project.org/package=BiocManager}}

\leavevmode\hypertarget{ref-biocpkgsub2020}{}%
\CSLLeftMargin{30. }
\CSLRightInline{{Package Submission} {[}Internet{]}. Bioconductor; 2020.
Available:
\url{https://www.bioconductor.org/developers/package-submission/}}

\leavevmode\hypertarget{ref-ropensci2020}{}%
\CSLLeftMargin{31. }
\CSLRightInline{{Transforming Science through Open Data and Software}
{[}Internet{]}. rOpenSci; 2020. Available: \url{https://ropensci.org/}}

\leavevmode\hypertarget{ref-ropenscidev2020}{}%
\CSLLeftMargin{32. }
\CSLRightInline{rOpenSci, Anderson B, Chamberlain S, Krystalli A, Mullen
L, Ram K, et al. {ropensci/dev\_guide: Fourth Release} {[}Internet{]}.
Zenodo; 2020.
doi:\href{https://doi.org/10.5281/zenodo.3749013}{10.5281/zenodo.3749013}}

\leavevmode\hypertarget{ref-drat}{}%
\CSLLeftMargin{33. }
\CSLRightInline{Carl Boettiger DE with contributions by, Fultz N, Gibb
S, Gillespie C, Górecki J, Jones M, et al. {drat: Drat R Archive
Template} {[}Internet{]}. 2020. Available:
\url{https://CRAN.R-project.org/package=drat}}

\leavevmode\hypertarget{ref-anderson2017}{}%
\CSLLeftMargin{34. }
\CSLRightInline{Anderson GB, Eddelbuettel D. {Hosting Data Packages via
drat: A Case Study with Hurricane Exposure Data}. R Journal. 2017;9. }

\leavevmode\hypertarget{ref-theussl2009}{}%
\CSLLeftMargin{35. }
\CSLRightInline{Theußl S, Zeileis A. {Collaborative Software Development
Using R-Forge. Special Invited Paper on "The Future of R"}. The R
Journal. The R Foundation for Statistical Computing; 2009;1: 9--14. }

\leavevmode\hypertarget{ref-lang2000}{}%
\CSLLeftMargin{36. }
\CSLRightInline{Lang DT. {The Omegahat Environment: New Possibilities
for Statistical Computing}. Journal of Computational and Graphical
Statistics. 2000;9: 423--451. }

\leavevmode\hypertarget{ref-mcelreath2020}{}%
\CSLLeftMargin{37. }
\CSLRightInline{McElreath R. {Statistical Rethinking: A Bayesian Course
with Examples in R and Stan}. CRC Press; 2020. }

\leavevmode\hypertarget{ref-decan2015}{}%
\CSLLeftMargin{38. }
\CSLRightInline{Decan A, Mens T, Claes M, Grosjean P. {On the
Development and Distribution of R packages: An Empirical Analysis of the
R Ecosystem}. {Proceedings of the 2015 European Conference on Software
Architecture Workshops}. 2015. pp. 1--6. }

\leavevmode\hypertarget{ref-devtools}{}%
\CSLLeftMargin{39. }
\CSLRightInline{Wickham H, Hester J, Chang W. {devtools: Tools to Make
Developing R Packages Easier} {[}Internet{]}. 2020. Available:
\url{https://CRAN.R-project.org/package=devtools}}

\leavevmode\hypertarget{ref-decan2016}{}%
\CSLLeftMargin{40. }
\CSLRightInline{Decan A, Mens T, Claes M, Grosjean P. {When GitHub meets
CRAN: An Analysis of Inter-Repository Package Dependency Problems}.
{2016 IEEE 23rd International Conference on Software Analysis,
Evolution, and Reengineering (SANER)}. IEEE; 2016. pp. 493--504. }

\leavevmode\hypertarget{ref-pkgdown}{}%
\CSLLeftMargin{41. }
\CSLRightInline{Wickham H, Hesselberth J. {pkgdown: Make Static HTML
Documentation for a Package} {[}Internet{]}. 2020. Available:
\url{https://CRAN.R-project.org/package=pkgdown}}

\leavevmode\hypertarget{ref-sf}{}%
\CSLLeftMargin{42. }
\CSLRightInline{Pebesma E. {Simple Features for R: Standardized Support
for Spatial Vector Data}. {The R Journal}. 2018;10: 439--446.
doi:\href{https://doi.org/10.32614/RJ-2018-009}{10.32614/RJ-2018-009}}

\leavevmode\hypertarget{ref-learnr}{}%
\CSLLeftMargin{43. }
\CSLRightInline{Schloerke B, Allaire J, Borges B. {learnr: Interactive
Tutorials for R} {[}Internet{]}. 2020. Available:
\url{https://CRAN.R-project.org/package=learnr}}

\leavevmode\hypertarget{ref-ushey2020}{}%
\CSLLeftMargin{44. }
\CSLRightInline{Ushey K. {RStudio 1.3 Preview: Integrated Tutorials}
{[}Internet{]}. RStudio; 2020. Available:
\url{https://blog.rstudio.com/2020/02/25/rstudio-1-3-integrated-tutorials/}}

\leavevmode\hypertarget{ref-pkgdownweb}{}%
\CSLLeftMargin{45. }
\CSLRightInline{Wickham H, Hesselberth J. {pkgdown} {[}Internet{]}.
RStudio; 2018. Available: \url{https://pkgdown.r-lib.org/}}

\leavevmode\hypertarget{ref-naniar}{}%
\CSLLeftMargin{46. }
\CSLRightInline{Tierney N, Cook D, McBain M, Fay C. {naniar: Data
Structures, Summaries, and Visualisations for Missing Data}
{[}Internet{]}. 2020. Available:
\url{https://CRAN.R-project.org/package=naniar}}

\leavevmode\hypertarget{ref-valr}{}%
\CSLLeftMargin{47. }
\CSLRightInline{Riemondy KA, Sheridan RM, Gillen A, Yu Y, Bennett CG,
Hesselberth JR. {valr: Reproducible Genome Interval Arithmetic in R}.
F1000Research. 2017;
doi:\href{https://doi.org/10.12688/f1000research.11997.1}{10.12688/f1000research.11997.1}}

\leavevmode\hypertarget{ref-TransPhylo}{}%
\CSLLeftMargin{48. }
\CSLRightInline{Didelot, Xavier, Fraser, Christophe, Gardy, Jennifer, et
al. {Genomic Infectious Disease Epidemiology in Partially Sampled and
Ongoing Outbreaks}. Molecular Biology and Evolution. 2017;34: 997--1007.
doi:\href{https://doi.org/10.1093/molbev/msw275}{10.1093/molbev/msw275}}

\leavevmode\hypertarget{ref-rdocumentation}{}%
\CSLLeftMargin{49. }
\CSLRightInline{{RDocumentation} {[}Internet{]}. DataCamp Inc.; 2020.
Available: \url{https://www.rdocumentation.org/}}

\leavevmode\hypertarget{ref-bookdown}{}%
\CSLLeftMargin{50. }
\CSLRightInline{Xie Y. {bookdown: Authoring Books and Technical
Documents with R Markdown} {[}Internet{]}. 2020. Available:
\url{https://CRAN.R-project.org/package=bookdown}}

\leavevmode\hypertarget{ref-xie2016}{}%
\CSLLeftMargin{51. }
\CSLRightInline{Xie Y. {bookdown: Authoring Books and Technical
Documents with R Markdown}. Chapman; Hall/CRC; 2016. }

\leavevmode\hypertarget{ref-mgcv}{}%
\CSLLeftMargin{52. }
\CSLRightInline{Wood SN. {Fast Stable Restricted Maximum Likelihood and
Marginal Likelihood Estimation of Semiparametric Generalized Linear
Models}. Journal of the Royal Statistical Society (B). 2011;73: 3--36. }

\leavevmode\hypertarget{ref-tidymodels}{}%
\CSLLeftMargin{53. }
\CSLRightInline{Kuhn M, Wickham H. {Tidymodels: A Collection of Packages
for Modeling and Machine Learning using Tidyverse Principles}
{[}Internet{]}. 2020. Available: \url{https://www.tidymodels.org}}

\leavevmode\hypertarget{ref-shiny}{}%
\CSLLeftMargin{54. }
\CSLRightInline{Chang W, Cheng J, Allaire J, Xie Y, McPherson J. {shiny:
Web Application Framework for R} {[}Internet{]}. 2020. Available:
\url{https://CRAN.R-project.org/package=shiny}}

\leavevmode\hypertarget{ref-htmlwidgets}{}%
\CSLLeftMargin{55. }
\CSLRightInline{Vaidyanathan R, Xie Y, Allaire J, Cheng J, Russell K.
{htmlwidgets: HTML Widgets for R} {[}Internet{]}. 2019. Available:
\url{https://CRAN.R-project.org/package=htmlwidgets}}

\leavevmode\hypertarget{ref-Rcpp}{}%
\CSLLeftMargin{56. }
\CSLRightInline{Eddelbuettel D, François R. {Rcpp: Seamless R and C++
Integration}. Journal of Statistical Software. 2011;40: 1--18.
doi:\href{https://doi.org/10.18637/jss.v040.i08}{10.18637/jss.v040.i08}}

\leavevmode\hypertarget{ref-caret}{}%
\CSLLeftMargin{57. }
\CSLRightInline{Kuhn M. {caret: Classification and Regression Training}
{[}Internet{]}. 2020. Available:
\url{https://CRAN.R-project.org/package=caret}}

\leavevmode\hypertarget{ref-datatable}{}%
\CSLLeftMargin{58. }
\CSLRightInline{Dowle M, Srinivasan A. {data.table: Extension of
data.frame} {[}Internet{]}. 2019. Available:
\url{https://CRAN.R-project.org/package=data.table}}

\leavevmode\hypertarget{ref-leaflet}{}%
\CSLLeftMargin{59. }
\CSLRightInline{Cheng J, Karambelkar B, Xie Y. {leaflet: Create
Interactive Web Maps with the JavaScript Leaflet Library}
{[}Internet{]}. 2019. Available:
\url{https://CRAN.R-project.org/package=leaflet}}

\leavevmode\hypertarget{ref-randomizr}{}%
\CSLLeftMargin{60. }
\CSLRightInline{Coppock A. {randomizr: Easy-to-Use Tools for Common
Forms of Random Assignment and Sampling} {[}Internet{]}. 2019.
Available: \url{https://CRAN.R-project.org/package=randomizr}}

\leavevmode\hypertarget{ref-stringr}{}%
\CSLLeftMargin{61. }
\CSLRightInline{Wickham H. {stringr: Simple, Consistent Wrappers for
Common String Operations} {[}Internet{]}. 2019. Available:
\url{https://CRAN.R-project.org/package=stringr}}

\leavevmode\hypertarget{ref-survminer}{}%
\CSLLeftMargin{62. }
\CSLRightInline{Kassambara A, Kosinski M, Biecek P. {survminer: Drawing
Survival Curves using ggplot2} {[}Internet{]}. 2020. Available:
\url{https://CRAN.R-project.org/package=survminer}}

\leavevmode\hypertarget{ref-Rmail2020}{}%
\CSLLeftMargin{63. }
\CSLRightInline{R Core Team. {Mailing Lists} {[}Internet{]}. The R
Foundation; 2020. Available: \url{https://www.r-project.org/mail.html}}

\leavevmode\hypertarget{ref-chase2020}{}%
\CSLLeftMargin{64. }
\CSLRightInline{Chase W. {Dataviz and the 20th Anniversary of R, an
Interview With Hadley Wickham} {[}Internet{]}. Medium; 2020. Available:
\url{https://medium.com/nightingale/dataviz-and-the-20th-anniversary-of-r-an-interview-with-hadley-wickham-ea245078fc8a}}

\leavevmode\hypertarget{ref-statnet}{}%
\CSLLeftMargin{65. }
\CSLRightInline{Handcock MS, Hunter DR, Butts CT, Goodreau SM, Krivitsky
PN, Morris M. {ergm: Fit, Simulate and Diagnose Exponential-Family
Models for Networks} {[}Internet{]}. The Statnet Project
(\url{http://www.statnet.org}); 2018. Available:
\url{https://CRAN.R-project.org/package=ergm}}

\leavevmode\hypertarget{ref-muenchen2012}{}%
\CSLLeftMargin{66. }
\CSLRightInline{Muenchen RA. {The Popularity of Data Analysis Software}.
URL http://r4statscom/popularity. 2012; }

\leavevmode\hypertarget{ref-wickham2019}{}%
\CSLLeftMargin{67. }
\CSLRightInline{Wickham H. {Advanced R}. CRC Press; 2019. }

\leavevmode\hypertarget{ref-eddelbuettel2013}{}%
\CSLLeftMargin{68. }
\CSLRightInline{Eddelbuettel D. {Seamless R and C++ Integration with
Rcpp}. Springer; 2013. }

\leavevmode\hypertarget{ref-githubdocs}{}%
\CSLLeftMargin{69. }
\CSLRightInline{{GitHub Docs} {[}Internet{]}. GitHub, Inc.; 2020.
Available: \url{https://docs.github.com/en}}

\leavevmode\hypertarget{ref-leek2015}{}%
\CSLLeftMargin{70. }
\CSLRightInline{Leek J. {How I Decide When to Trust an R Package}
{[}Internet{]}. 2015. Available:
\url{https://simplystatistics.org/2015/11/06/how-i-decide-when-to-trust-an-r-package/}}

\leavevmode\hypertarget{ref-vannoorenberghe2017}{}%
\CSLLeftMargin{71. }
\CSLRightInline{Vannoorenberghe L. {RDocumentation: Scoring and Ranking}
{[}Internet{]}. DataCamp; 2017. Available:
\url{https://www.datacamp.com/community/blog/rdocumentation-ranking-scoring}}

\leavevmode\hypertarget{ref-packagemetrics}{}%
\CSLLeftMargin{72. }
\CSLRightInline{Firke S, Krouse B, Grand E, Shepherd L, Ampeh W, Frick
H. {packagemetrics: A Package for Helping You Choose Which Package to
Use} {[}Internet{]}. 2017. Available:
\url{https://github.com/ropenscilabs/packagemetrics}}

\leavevmode\hypertarget{ref-packagemetricsblog}{}%
\CSLLeftMargin{73. }
\CSLRightInline{Becca Krouse HF Erin Grand. {packagemetrics - Helping
you Choose a Package Since runconf17} {[}Internet{]}. rOpenSci; 2017.
Available: \url{https://ropensci.org/blog/2017/06/27/packagemetrics/}}

\leavevmode\hypertarget{ref-tidyverse}{}%
\CSLLeftMargin{74. }
\CSLRightInline{Wickham H, Averick M, Bryan J, Chang W, McGowan LD,
François R, et al. {Welcome to the Tidyverse}. Journal of Open Source
Software. 2019;4: 1686.
doi:\href{https://doi.org/10.21105/joss.01686}{10.21105/joss.01686}}

\leavevmode\hypertarget{ref-smith2014}{}%
\CSLLeftMargin{75. }
\CSLRightInline{Smith A. {Improving GitHub for Science} {[}Internet{]}.
GitHub, Inc.; 2014. Available:
\url{https://github.blog/2014-05-14-improving-github-for-science/}}

\leavevmode\hypertarget{ref-flowCore}{}%
\CSLLeftMargin{76. }
\CSLRightInline{Ellis B, Haaland P, Hahne F, Le Meur N, Gopalakrishnan
N, Spidlen J, et al. {flowCore: Basic Structures for Flow Cytometry
Data}. 2019. }

\leavevmode\hypertarget{ref-hahne2009}{}%
\CSLLeftMargin{77. }
\CSLRightInline{Hahne F, LeMeur N, Brinkman RR, Ellis B, Haaland P,
Sarkar D, et al. {flowCore: A Bioconductor Package for High Throughput
Flow Cytometry}. BMC bioinformatics. BioMed Central; 2009;10: 1--8. }

\leavevmode\hypertarget{ref-taschuk2017}{}%
\CSLLeftMargin{78. }
\CSLRightInline{Taschuk M, Wilson G. {Ten Simple Rules for Making
Research Software more Robust}. PLoS Computational Biology. 2017;13. }

\leavevmode\hypertarget{ref-perez2016}{}%
\CSLLeftMargin{79. }
\CSLRightInline{Perez-Riverol Y, Gatto L, Wang R, Sachsenberg T,
Uszkoreit J, Veiga Leprevost F da, et al. {Ten Simple Rules for Taking
Advantage of Git and GitHub}. PLoS Computational Biology. 2016;12. }

\leavevmode\hypertarget{ref-hester2020}{}%
\CSLLeftMargin{80. }
\CSLRightInline{Hester J. {How Does covr Work Anyway?} {[}Internet{]}.
The R Foundation; 2020. Available:
\url{https://cran.r-project.org/web/packages/covr/vignettes/how_it_works.html}}

\leavevmode\hypertarget{ref-wickham2011}{}%
\CSLLeftMargin{81. }
\CSLRightInline{Wickham H. {testthat: Get Started with Testing}. The R
Journal. 2011;3: 5--10. Available:
\url{https://journal.r-project.org/archive/2011-1/RJournal_2011-1_Wickham.pdf}}

\leavevmode\hypertarget{ref-covr}{}%
\CSLLeftMargin{82. }
\CSLRightInline{Hester J. {covr: Test Coverage for Packages}
{[}Internet{]}. 2020. Available:
\url{https://CRAN.R-project.org/package=covr}}

\leavevmode\hypertarget{ref-barts2018}{}%
\CSLLeftMargin{83. }
\CSLRightInline{Barts C. {How to use GitHub Badges to Stop Feeling Like
a Noob} {[}Internet{]}. freeCodeCamp; 2018. Available:
\url{https://www.freecodecamp.org/news/how-to-use-badges-to-stop-feeling-like-a-noob-d4e6600d37d2/}}

\leavevmode\hypertarget{ref-wilson2006}{}%
\CSLLeftMargin{84. }
\CSLRightInline{Wilson GV. {Where's the Real Bottleneck in Scientific
Computing?} American Scientist. 2006;94: 5. }

\leavevmode\hypertarget{ref-bryan2018}{}%
\CSLLeftMargin{85. }
\CSLRightInline{Bryan J. {Excuse Me, Do You Have a Moment to Talk about
Version Control?} The American Statistician. Taylor \& Francis; 2018;72:
20--27. }

\leavevmode\hypertarget{ref-perkel2016}{}%
\CSLLeftMargin{86. }
\CSLRightInline{Perkel J. {Democratic Databases: Science on GitHub}.
Nature News. 2016;538: 127. }

\leavevmode\hypertarget{ref-peng2017}{}%
\CSLLeftMargin{87. }
\CSLRightInline{Roger D. Peng Sean Kross, Anderson B. {Mastering
Software Development in R}. https://rdpeng.github.io/RProgDA/; 2017. }

\leavevmode\hypertarget{ref-broom}{}%
\CSLLeftMargin{88. }
\CSLRightInline{Robinson D, Hayes A. {broom: Convert Statistical
Analysis Objects into Tidy Tibbles} {[}Internet{]}. 2020. Available:
\url{https://CRAN.R-project.org/package=broom}}

\leavevmode\hypertarget{ref-biobroom}{}%
\CSLLeftMargin{89. }
\CSLRightInline{Andrew J. Bass SL David G. Robinson. {biobroom: Turn
Bioconductor Objects into Tidy Data Frames} {[}Internet{]}. 2020.
Available: \url{https://github.com/StoreyLab/biobroom}}

\leavevmode\hypertarget{ref-parker2013}{}%
\CSLLeftMargin{90. }
\CSLRightInline{Parker H. {Personal R Packages} {[}Internet{]}. 2013.
Available:
\url{https://hilaryparker.com/2013/04/03/personal-r-packages/}}

\leavevmode\hypertarget{ref-Hmisc}{}%
\CSLLeftMargin{91. }
\CSLRightInline{Harrell Jr FE, Charles Dupont with contributions from,
others. many. {Hmisc: Harrell Miscellaneous} {[}Internet{]}. 2020.
Available: \url{https://CRAN.R-project.org/package=Hmisc}}

\leavevmode\hypertarget{ref-broman}{}%
\CSLLeftMargin{92. }
\CSLRightInline{Broman KW. {broman: Karl Broman's R Code}
{[}Internet{]}. 2020. Available:
\url{https://CRAN.R-project.org/package=broman}}

\leavevmode\hypertarget{ref-robinson2014}{}%
\CSLLeftMargin{93. }
\CSLRightInline{Robinson D. {broom: An R package for Converting
Statistical Analysis Objects into Tidy Data Frames}. arXiv preprint
arXiv:14123565. 2014; }

\leavevmode\hypertarget{ref-usethis}{}%
\CSLLeftMargin{94. }
\CSLRightInline{Wickham H, Bryan J. {usethis: Automate Package and
Project Setup} {[}Internet{]}. 2019. Available:
\url{https://CRAN.R-project.org/package=usethis}}

\leavevmode\hypertarget{ref-roxygen2}{}%
\CSLLeftMargin{95. }
\CSLRightInline{Wickham H, Danenberg P, Csárdi G, Eugster M. {roxygen2:
In-Line Documentation for R} {[}Internet{]}. 2020. Available:
\url{https://CRAN.R-project.org/package=roxygen2}}

\leavevmode\hypertarget{ref-rlang}{}%
\CSLLeftMargin{96. }
\CSLRightInline{Henry L, Wickham H. {rlang: Functions for Base Types and
Core R and Tidyverse Features} {[}Internet{]}. 2020. Available:
\url{https://CRAN.R-project.org/package=rlang}}

\leavevmode\hypertarget{ref-Rcore2020}{}%
\CSLLeftMargin{97. }
\CSLRightInline{R Core Team. {Writing R Extensions} {[}Internet{]}. The
R Foundation; 2020. Available:
\url{https://cran.r-project.org/doc/manuals/R-exts.html}}

\end{CSLReferences}

\nolinenumbers


\end{document}
